% mss.tex
% Evidence: ../PLAN.md, ../survey/S0--S9, ../data/README.md, and lean-pokemon (branch master,
% commit ce0de3b, read from the read-only checkout ../lean-pokemon-ci). Every claim below
% should trace to one of these. Lean statements are copied verbatim from the source; do not
% edit them by hand -- re-copy them. The revision history of this manuscript is in git.

\documentclass[11pt]{article}

\usepackage[utf8]{inputenc}
\usepackage[T1]{fontenc}
\usepackage{lmodern}
\usepackage[margin=1in]{geometry}
\usepackage{amsmath,amssymb,amsthm}
\usepackage{booktabs}
\usepackage{tabularx}
\usepackage{xltabular}
\usepackage{xcolor}
\usepackage{listings}
\usepackage{microtype}
\usepackage{hyperref}

\newcommand{\pkmn}{Pok\'emon}
\newcommand{\lean}[1]{\texttt{#1}}
\newcommand{\todo}[1]{\textcolor{red!70!black}{[\textsc{todo}: #1]}}
\newcommand{\ev}[1]{\textsf{[#1]}}

\theoremstyle{plain}
\newtheorem{theorem}{Theorem}
\newtheorem{lemma}[theorem]{Lemma}
\newtheorem{corollary}[theorem]{Corollary}
\newtheorem{conjecture}[theorem]{Conjecture}
\theoremstyle{definition}
\newtheorem{definition}[theorem]{Definition}
\theoremstyle{remark}
\newtheorem{remark}[theorem]{Remark}

\lstdefinelanguage{lean4}{
  morekeywords={theorem,def,inductive,structure,where,deriving,namespace,end,fun,let,
    import,Prop,Type,Nat,Bool,List,Option,some,none,if,then,else,match,with},
  sensitive=true,
  morecomment=[l]{--},
  morecomment=[s]{/-}{-/},
  morestring=[b]",
}
\lstset{
  language=lean4,
  basicstyle=\ttfamily\upshape\footnotesize,
  keywordstyle=\bfseries,
  commentstyle=\itshape\color{black!60},
  columns=fullflexible,
  keepspaces=true,
  breaklines=true,
  frame=single,
  framerule=0.3pt,
  xleftmargin=1em,
  xrightmargin=1em,
  aboveskip=0.8em,
  belowskip=0.8em,
  literate=
    {→}{{$\to$}}1
    {∀}{{$\forall$}}1
    {∃}{{$\exists$}}1
    {∧}{{$\wedge$}}1
    {∨}{{$\vee$}}1
    {¬}{{$\neg$}}1
    {≤}{{$\le$}}1
    {≠}{{$\ne$}}1
    {↔}{{$\leftrightarrow$}}1
    {∈}{{$\in$}}1
    {≥}{{$\ge$}}1
    {·}{{$\cdot$}}1
    {₀}{{$_0$}}1
    {₁}{{$_1$}}1
    {₂}{{$_2$}}1
    {₃}{{$_3$}}1
    {é}{{\'e}}1
    {×}{{$\times$}}1
    {⟨}{{$\langle$}}1
    {⟩}{{$\rangle$}}1,
}

\title{Finite or Universal:\\ Mechanized Results for a Model of\\ the \pkmn{} Trading Card Game}
\author{N. E. Davis (Sigilante)}
\date{}

\begin{document}
\maketitle

\begin{abstract}
Magic, Yu-Gi-Oh!\ and Riftbound embed universal computation under stated assumptions about
play. We ask whether the \pkmn{} Trading Card Game does, for a Lean~4 model of its rules, not
the printed game. The model has Special Conditions, Checkup, rulebook Knock Out timing and
damage-reading effects. We prove three results, machine-checked under an axiom audit. First,
for tame decks (no card adds a damage reduction or heals by the damage it does), from a
well-formed start without reductions, a clamp capping damage, Poison strength and damage
bonuses has finitely many reachable values and is a bisimulation for moves. Second, for every
well-formed start whose cards put no damage bonus or reduction inside a repeat and give
rest-of-game bonuses or reductions in attacks only through once-per-game GX attacks (damage
reduction, healing by damage done, carried modifiers and Trainer-card modifiers all allowed),
the reachable states fall into finitely many classes of a two-way bisimulation that keeps every
modifier exact. The bound uses the model's lack of card recycling, repeated attacks and extra
turns, the first two of which today's Standard cards have. Third, without the GX restriction, with the further generality of
Section~\ref{sec:counter-generality}, decks exist that reach a halt marker by some sequence of legal
moves exactly when a given two-counter machine halts; with Minsky's theorem (cited, not
formalised) this reachability is undecidable. Nothing follows for the printed game, which we
conjecture is not Turing complete. A non-exhaustive survey of 20,776 English card records found
no unbounded store that card text reads beyond a yes-or-no test, except possibly via
card-recycling loops.
\end{abstract}

\noindent{\small\copyright{} 2026 N.~E.~Davis (Sigilante). This work is licensed under the Creative Commons Attribution 4.0 International licence (CC~BY~4.0), \url{https://creativecommons.org/licenses/by/4.0/}.}


\tableofcontents

%======================================================================
\section{Introduction}\label{sec:intro}
%======================================================================

\subsection{The question and its precedents}

A game is \emph{Turing complete}, informally, if a position of the game can be built whose
continued play carries out an arbitrary computation, so that some question about the play
(who wins, whether a position is reachable) is as hard as the halting problem. The question
is precise only under assumptions about play: which choices the players make, whether
randomness is allowed, and which feature of the game stands in for unbounded memory. This
paper asks the question of the \pkmn{} Trading Card Game (TCG). Our search of the literature
found no published claim, proof or refutation of its Turing completeness
(Section~\ref{sec:related}); the only mention we found is a sentence in the Yu-Gi-Oh!\ paper
that asks whether its result ``holds for other card games (like Pokemon TGC)'' \cite{ygo}.

Churchill, Biderman and Herrick showed, in a peer-reviewed paper, that determining the outcome
of a game of Magic: The Gathering in which all remaining moves are forced is undecidable
\cite{cbh-mtg}. Two features of Magic carry their construction: \emph{forced play}, since
stacked libraries and mandatory triggered abilities remove every choice from both players; and
\emph{unbounded memory}, since the tape is a line of creature tokens and the game creates a
fresh token whenever the head reaches the end of the tape. Two preprints obtain universality
without forced play, and they show that tokens are not necessary. Nicolosi, Pisciotta and
Bresolin \cite{ygo} show that deciding whether a given computable strategy wins from a given
Yu-Gi-Oh!\ position is undecidable (indeed $\Pi^1_1$-complete when the second player may
choose natural numbers); their memory is a single uncapped number, the count of spell counters
on one card. Dalaklis and Fields \cite{riftbound} embed Turing machines in Riftbound under
``game-state-directed'' strategies, using tokens produced by an unbounded energy loop within
a single turn. All three results turn on a quantity with no cap that card text can read.

The \pkmn{} TCG offers neither feature in an evident form. As far as our survey found, no card
creates a game object or brings a card in from outside the game, so the 120 cards of two
60-card decks are all the objects there are. The numeric quantity closest to a register, the
number of damage counters on a \pkmn{}, is capped by the Knock Out rule. And no card forces a
player to attack, while every main-phase action is optional, so the fully forced play of
\cite{cbh-mtg} has no evident analogue (Section~\ref{sec:control}). The question for \pkmn{} is
therefore whether any uncapped quantity exists and can be read by card text.

\subsection{Results}\label{sec:results-intro}

We answer the question for a model of the rules, written in Lean~4 \cite{lean-pokemon}, and
report evidence about how far the answer reaches the printed game. The target is the game as it
stands today, under its current rules, not earlier versions of them. The program is to index
every printed card eventually and to prove results over today's Standard format first; \pkmn{}
made of several cards (V-UNION, LEGEND) are out of scope. The model has the rulebook's turn structure, typed attack costs, Weakness and
Resistance, the rulebook's damage calculation made once per attack, Special Conditions, a
between-turns Checkup phase, Knock Out timing, markers kept as flags, lingering damage
modifiers, and a language of card effects that can test the damage on either Active \pkmn{}
(Section~\ref{sec:model}). There is one model, and we prove three results about it: finiteness
for tame decks, finiteness for faithful decks, and a two-counter reduction for decks that drop
one clause of faithfulness and that no printed card supplies. A result labelled
\emph{machine-checked} is a theorem of the development at commit \texttt{ce0de3b}. The
development builds in continuous integration, and the build runs an audit that fails if any
theorem depends on an axiom other than \lean{propext} and \lean{Quot.sound}; a build at that
commit reports 458 theorems with source declarations and 1,122 generated ones, all passing
(Appendix~\ref{app:audit}).

\paragraph{Finiteness for tame decks (machine-checked).} Call a deck \emph{tame} if no card in it adds a
damage reduction or uses \lean{drain}, a heal by the damage done. The \emph{clamp} is the map
on states that caps damage and Poison strength (the extra damage counters that Poison places at
each Checkup) on the two Active \pkmn{} at the printed HP, and every stored damage bonus and
reduction at $2k$, where the \emph{cap} $k$ is a number computed from the cards that exceeds
every printed HP, every printed Resistance and every damage constant that card text tests. The
factor 2 is there because Resistance is subtracted after the bonus is added. A start is
\emph{well formed} if both
players have Prize cards left and no Benched \pkmn{} is at its HP or has a Special Condition.
From a well-formed start whose cards are tame and in which no reduction is in effect, the clamps of
the reachable states lie in one finite list, and states with the same clamp have the same moves
up to the clamp (Theorem~\ref{thm:finite}). So the clamped game is a finite transition system
with the same moves. The unclamped reachable space is infinite (Theorem~\ref{thm:infinite}),
and tameness is needed: for some decks that are not tame, no cap makes the clamp a simulation
(Theorem~\ref{thm:breaks}). Effects read damage, and inside one attack the clamp is not a
simulation (Theorem~\ref{thm:notsim}); for tame decks it is one between moves, because of the
model's Knock Out timing.

\paragraph{Finiteness for faithful decks (machine-checked).} A second clamp, the \emph{exact-store
clamp}, caps only damage and Poison strength on the two Active \pkmn{} and keeps every lingering
modifier exact. For any two states with the same exact-store clamp, every move from either is
matched by a move from the other, with no hypothesis on the deck (Theorem~\ref{thm:xbisim}).
Call a start \emph{faithful} if no card puts a damage bonus or reduction inside a repeat
(\lean{iter}) and attacks write rest-of-game bonuses or reductions only if they are GX attacks,
which the model allows once per player per game. From a well-formed faithful start, the
reachable states fall into finitely many classes of the exact-store clamp, with the two-way
bisimulation (Theorem~\ref{thm:xfinite}). Faithful decks may use damage reduction,
\lean{drain}, carried modifiers and modifiers on Item and Supporter cards. The bound on the
modifiers rests on three absences in the model, each of which exists on printed cards: no effect
returns a card from a discard pile, one attack is made per turn, and there are no extra turns.
Today's Standard pool has the first two, so the bound does not yet cover it
(Section~\ref{sec:standard}).
Neither class contains the other, and each theorem applies to its own class
(Section~\ref{sec:faithful}).

\paragraph{The universal side (machine-checked reduction; cited undecidability).} For decks
that no printed card supplies, the same model simulates two-counter machines. For every
two-counter machine $M$ there is a well-formed start state from which a halt marker (a marker
that the construction places when $M$ halts) is reachable by some sequence of legal moves of the
two players if and only if $M$ halts (Theorem~\ref{thm:universal}). The ``if'' direction is
cooperative reachability; the ``only if'' direction holds over all play. The construction uses rest-of-game damage modifiers granted by
an ordinary attack every turn and accumulating without limit, together with the further
generality listed in Section~\ref{sec:counter-generality}. Its start states satisfy the first
clause of faithfulness and violate the second, since their attacks are not GX attacks
(Theorem~\ref{thm:sharp}); so the GX clause cannot be dropped from the faithful theorem. This is
a sharpness result for that clause, not a classification of decks: nothing is proved about
other decks that are not faithful, and some of them (tame decks that violate the GX clause) are
finite. With Minsky's
theorem that halting of two-counter machines is undecidable \cite{minsky}, reachability of the
halt marker over the constructed start states is undecidable; the reduction is
machine-checked, and the appeal to Minsky's theorem is not.

\paragraph{What remains open.} The faithful class leaves out printed behaviours that the model
lacks: card recycling, repeated attacks and extra turns, which would reopen the bound, and
copy attacks, static modifiers and modifiers from Tools, Abilities or Special Energy,
reductions restricted by type, and locks on GX attacks. The open case is concrete: today's
Standard decks are not tame, since they carry lingering damage reductions and \lean{drain}
attacks, so the faithful theorem is the one that bears on them, and it needs card recycling and
repeated attacks in the model before it can. Whether finiteness survives that extension, for the
Standard pool, is open, and so is a second route to universality that would need a replayable
card writing two coupled modifiers (Section~\ref{sec:linger-next}).

\paragraph{The printed card pool (surveyed).} A regular-expression scan and a template census
of 20,776 English card records, followed by reading every sentence in the risk classes, found
no unbounded store that card text reads beyond a yes-or-no test (Section~\ref{sec:survey}).
The survey is not exhaustive: it covers English text only, it finds only the wordings its
patterns describe, and it leaves one case open, since a card-recycling loop that replays a
damage modifier within one turn would make that turn's modifier unbounded. A profile of today's
Standard pool found no rest-of-game modifiers, so the universal construction has no Standard
instance (Section~\ref{sec:standard}).

\paragraph{Conjectures.} We conjecture that the printed \pkmn{} TCG is not Turing complete
(Conjecture~\ref{conj:main}), and, with less evidence, that unbounded HP alone would make it
universal (Section~\ref{sec:lattice}).

\subsection{Status of claims}\label{sec:status}

Table~\ref{tab:status} lists the headline claims with the kind of evidence behind each;
Appendix~\ref{app:claims} lists every claim. \emph{Machine-checked} means a Lean theorem at
commit \texttt{ce0de3b} \cite{lean-pokemon} that passes the axiom audit of
Appendix~\ref{app:audit}. \emph{Argued} means a proof by hand, not mechanized.
\emph{Surveyed} means a systematic but non-exhaustive search of rules documents or card text.
\emph{Conjectured} means that no proof is offered, and \emph{open} that we state a question and
take no side.

\begin{table}[ht]
\centering
\footnotesize
\begin{tabularx}{\textwidth}{@{}>{\raggedright\arraybackslash}X >{\raggedright\arraybackslash}p{2.1cm} >{\raggedright\arraybackslash}p{3.0cm} l@{}}
\toprule
Claim & Decks & Evidence & Where \\
\midrule
Legal play permutes the cards in the game & all & machine-checked & Thm.~\ref{thm:perm} \\
An attack ends the turn; with consistent turn flags the first-turn bans hold & all & machine-checked & Thm.~\ref{thm:turn} \\
The unclamped reachable set is infinite, even without a pending Knock Out & one witness & machine-checked & Thm.~\ref{thm:infinite} \\
No move and no step of an Item or Supporter acts on a \pkmn{} at or above its HP, except to Knock it Out & all & machine-checked & Thms.~\ref{thm:untouched}, \ref{thm:guarded} \\
The model's Knock Out timing inside Item and Supporter effects is the Advanced Player's Rulebook's, up to two caveats; no English rulebook states it & all & argued & Sec.~\ref{sec:departures} \\
The clamp is not a simulation inside one attack's effect & one witness per cap & machine-checked & Thm.~\ref{thm:notsim} \\
Clamps of reachable states lie in one finite list, and the clamp (stores capped at $2k$) is a bisimulation for moves & tame & machine-checked & Thm.~\ref{thm:finite} \\
No cap makes the clamp a simulation & one witness per cap, with a reduction or \lean{drain} & machine-checked & Thm.~\ref{thm:breaks} \\
States with the same exact-store clamp have the same moves, both ways & all & machine-checked & Thm.~\ref{thm:xbisim} \\
Reachable states lie in finitely many exact-store classes, with the two-way bisimulation & faithful & machine-checked & Thm.~\ref{thm:xfinite} \\
No fixed start simulates a two-counter machine with infinitely many reachable configurations & tame or faithful & argued & Cor.~\ref{cor:counter} \\
A halt marker is reachable by some sequence of legal moves iff a given two-counter machine halts & built for each machine; loop-free, not GX-only & machine-checked & Thms.~\ref{thm:universal}, \ref{thm:sharp} \\
Reachability of the halt marker over those starts is undecidable & built for each machine & machine-checked reduction; Minsky's theorem cited & Sec.~\ref{sec:counter-what} \\
Finiteness for faithful decks once card recycling and repeated attacks are modelled & today's Standard pool & open & Sec.~\ref{sec:linger-next} \\
A replayable card that writes two coupled modifiers gives universality without rest-of-game modifiers & with recycling or repeated attacks & argued; not checked by a second reader & Sec.~\ref{sec:linger-next} \\
No unbounded store that card text reads beyond a yes-or-no test, except possibly through card-recycling loops & English card pool & surveyed & Sec.~\ref{sec:survey} \\
Today's Standard pool has card recycling, attacking twice, lingering reductions and \lean{drain}, but no rest-of-game modifiers, GX attacks or extra turns & Standard (by regulation mark) & surveyed & Sec.~\ref{sec:standard} \\
The printed game is not Turing complete & English card pool & conjectured & Conj.~\ref{conj:main} \\
Unbounded HP would make the game universal & printed, HP unbounded & conjectured & Sec.~\ref{sec:lattice} \\
\bottomrule
\end{tabularx}
\caption{Headline claims. The full list is in Appendix~\ref{app:claims}.}\label{tab:status}
\end{table}

\paragraph{Plan of the paper.} Section~\ref{sec:game} explains the game for readers who do not
play it, and Section~\ref{sec:model} presents the Lean model and its departures from the
rulebook. Section~\ref{sec:finite-side} proves the two finiteness theorems and
Section~\ref{sec:linger} the two-counter reduction. Section~\ref{sec:printed} turns to the
printed game: the survey, rules coverage and today's Standard pool, where memory could hide,
what remains open, and the conjectures. Section~\ref{sec:related} discusses related
work, and the appendices hold the full departure list, the state inventory, survey statistics,
the Lean audit, the full claims table and future work.

%======================================================================
\section{The game for non-players}\label{sec:game}
%======================================================================

The reference is the game's current rules. Three English rulebooks bear on them. The one on
hand is the 2023 edition, the Scarlet \& Violet---151 rulebook \cite{rb-151}, which we read in
full and cite by page. The 2024 rulebook titled ``\pkmn{} Trading Card Game Rules''
\cite{rb-ssp} (we cite it as the Scarlet \& Violet---Surging Sparks rulebook, an attribution
that rests on its file name, \texttt{ssp\_rulebook\_en.pdf}; the text does not name the set)
was read from a third-party copy, and a 2026 post on rule changes \cite{pp26} only through a
summarising fetch. The current rulebook has not been read directly: a rule changed since 2024
may be missing from the model, and the 2026 changes we cite are unverified against a rulebook.
On the rules that the model uses and that we compared (the damage calculation, Checkup, Special
Conditions, the first-turn rules and the Stadium rule), the 2023 and 2024 rulebooks agree. We
also use the
Advanced Player's Rulebook Ver.~3.4 of \pkmn{} Asia \cite{apr}, tournament documents
\cite{trh26,pg26,tr2013}, and two rulings \cite{compendium}. Our rules audit read the 2024
rulebook and these documents in full; Appendix~\ref{app:inventory} lists the state components
it found, with a bound for each.

\paragraph{Cards and zones.} Each player has a deck of exactly 60 cards \cite{rb-ssp}, so a
game has $N = 120$ cards. A \emph{\pkmn{}} card has a printed HP (hit points); a \emph{Basic}
\pkmn{} is put into play directly, and an evolved one is played on top of the \pkmn{} it
evolves from. Each player has one \emph{Active} \pkmn{}, in the \emph{Active Spot}, which
attacks and is attacked, and a \emph{Bench} of up to five more. \emph{Energy} cards are
attached to \pkmn{}, one per turn, and pay for attacks. \emph{Trainer} cards are of four kinds:
an \emph{Item} may be played any number of times in a turn, a \emph{Supporter} once per turn, a
\emph{Tool} is attached to a \pkmn{}, and a \emph{Stadium} stays in play and affects both
players. At set-up each player puts six cards aside face down as \emph{Prize cards}. The other
zones are the deck (ordered), the hand, the discard pile and the Lost Zone. \emph{Card
conservation} holds: no rule creates or destroys a physical card, and effects only move cards
among zones, attach them or stack them.

\paragraph{Damage and Knock Outs.} Damage is placed as \emph{damage counters}, one for each
10~HP. A \pkmn{} whose damage counters reach its HP is \emph{Knocked Out}: it and its attached
cards go to the discard pile, and its opponent takes one Prize card. A player wins by taking
the last Prize card, or when the opponent has no \pkmn{} in play or cannot draw at the start of
a turn. The rulebook states that ``if a \pkmn{} has total damage at least equal to its HP
\dots\ it is Knocked Out'' \cite{rb-ssp}, and it checks for Knock Outs after an attack
(``check to see if any \pkmn{} were Knocked Out by the attack'', \cite[p.~14]{rb-151}) and after
Checkup (``After both players have gone through these checks'', \cite[p.~15]{rb-151}). It
states no check at any other time. The Advanced Player's Rulebook of \pkmn{} Asia, a document of
another region, adds one: ``At the end of an attack'', ``At the end of \pkmn{} checkup'', and
``Anytime a \pkmn{}'s remaining HP becomes zero, in addition to the above'' \cite{apr}. We
found no English rule for the third clause, and the model's Knock Outs at other times rest on
that document alone. Two rulings
confirm that damage may exceed HP transiently: inside one attack (``you can put up to 9 damage
counters on it even if it doesn't have that much HP remaining''), and inside one Checkup, where
the Knock Out waits until the between-turns step is over \cite{compendium}. So a \pkmn{} has
more damage counters than its HP only within one attack or one Checkup, and at the next check
it is Knocked Out.

\paragraph{Turns.} A turn begins with a draw. The player may then, in any order and with the
limits stated, put Basic \pkmn{} on the Bench, attach one Energy, evolve, \emph{retreat} the
Active \pkmn{} (swap it with a Benched one, paying Energy) once, play Items and one Supporter,
and use \emph{Abilities}, card text that is not an attack. The turn ends with an optional
attack: ``Once you attack, your turn is over'' \cite[p.~13]{rb-151}. On the first turn of the
game the starting player may neither attack nor play a Supporter, and neither player may evolve
on their own first turn \cite[pp.~11--13]{rb-151}. An attack's cost is a list of Energy
symbols: a typed symbol needs Energy of that type, and a Colorless symbol any Energy. Its damage
is computed once, in a fixed order \cite[p.~20]{rb-151}: the base damage (stop if it is 0), then
the attacker's bonuses, then the defender's \emph{Weakness} and \emph{Resistance} to the
attacker's type (on today's Standard cards Weakness is $\times 2$ and Resistance $-30$, once
$-20$), then the defender's reductions, each subtraction stopping at 0. An effect that places
damage counters directly does not do damage and skips the calculation. A few powers, the GX and VSTAR attacks and abilities, may be used only
once per game. Many printed effects \emph{linger}: they last ``during your opponent's next turn'' or,
rarely, for the rest of the game.

\paragraph{Special Conditions and Checkup.} An Active \pkmn{} may be Asleep, Confused or
Paralyzed (these exclude one another), and independently Poisoned and Burned; these are the five
\emph{Special Conditions}. A new Poison or Burn replaces the old one \cite[p.~16]{rb-151}. A
\pkmn{} loses its Special Conditions when it moves to the Bench or evolves.
\emph{\pkmn{} Checkup} runs between turns: Poison and Burn place damage counters, an Asleep
\pkmn{} flips a coin to wake, Paralysis ends after its owner's turn, and then Knock Outs are
checked.

\paragraph{Time and memory.} Items, Benching and evolution have no per-turn limit
\cite{rb-ssp}, so a turn can contain unboundedly many actions, and the game has no turn limit;
as far as our rules audit found, no rule forbids infinite loops
\cite{rb-ssp,apr,trh26,pg26,tr2013}. These unbounded quantities are \emph{time}, not memory.
The rules alone bound the state between Knock Out checks, provided four things hold: card
conservation, damage bounded by HP through the Knock Out check, HP bounded, and all other
memory being finitely many flags. Whether HP, markers and lingering effects are bounded is a
question about card text, which Section~\ref{sec:survey} asks.

%======================================================================
\section{The model}\label{sec:model}
%======================================================================

The model is the Lean~4 development \cite{lean-pokemon}, branch \texttt{master}, commit
\texttt{ce0de3b}, built with core Lean v4.36.0 and no libraries. HP and damage are measured in
damage counters. All Lean text in this paper is copied from that commit; docstrings and
\lean{deriving} clauses are omitted unless shown.

\subsection{State}

A \pkmn{} in play (\lean{InPlay}) is a record of its top card, the cards beneath it, attached
Energy and Tools, a number of damage counters, a rotation condition \lean{rot} (none, Asleep,
Confused, Paralyzed), two flags \lean{poisoned} and \lean{burned}, a number \lean{poisonUp} of
extra Poison damage counters (the \emph{Poison strength}: a Poisoned \pkmn{} takes
$1 + {}$\lean{poisonUp} damage counters at each Checkup), a flag \lean{fresh} for ``put into
play or evolved this turn'', and its markers \lean{marks} (Section~\ref{sec:markers}). A
\pkmn{} card records its HP, type, attacks, Retreat Cost, an optional Weakness (a type) and an
optional Resistance (a type and an amount). Energy cards are Basic Energy of the nine types
(Grass, Fire, Water, Lightning, Psychic, Fighting, Darkness, Metal, Fairy); there is no Special
Energy. A player's state holds the deck (head is the top), hand, discard pile, Prizes, Lost
Zone, an optional Active \pkmn{}, a Bench, three once-per-turn flags, two once-per-game flags
(\lean{gxUsed}, which records that the player has used a GX attack, and a VSTAR flag, which no
rule uses), a store \lean{mods} of lingering modifiers (Section~\ref{sec:linger-mech}), and a
flag \lean{hadTurn}, set when the player's first turn ends. The game state \lean{GameState}
(Appendix~\ref{app:defs}) holds the two players' states, \lean{cur} for the player to move and
\lean{opp} for the other, and six global components: whether the player to move went first,
whether this is the first turn of the game, whether the player to move has attacked this turn,
whether the game is between turns (\lean{between}), the winner, if any, and one Stadium slot
shared by both players, which only holds a card of the start state, since no move plays a
Stadium. The model does not record the turn number, since card text reads the turn number only
as ``first turn'', which \lean{firstTurn} and \lean{hadTurn} record. \lean{GameState.allCards}
lists every physical card in every zone of both players and in the Stadium slot.

\subsection{The effect language and \lean{Run}}\label{sec:effects}

Card text is represented by a small language. A \emph{test} reads one Active \pkmn{}, the
player's own (\lean{own}) or the opponent's (\lean{opp}):

\begin{lstlisting}
inductive Test where
  | atLeast (w : Whose) (n : Nat)   -- it has at least n damage counters
  | exactly (w : Whose) (n : Nat)   -- it has exactly n damage counters
  | cond (w : Whose) (c : Condition) -- it is affected by the Special Condition c
  | marked (w : Whose) (m : Marker)  -- it has the marker m
  | isType (w : Whose) (t : EType)   -- its type is t (top card's, or set by Coloring)
\end{lstlisting}
\noindent ``If it has no damage counters'' is the \lean{else} branch of \lean{atLeast w 1};
Mega Absol ex's exact-6 test is \lean{exactly opp 6}. An \emph{amount} is the base damage that
an attack's text computes, in damage counters, before any modifier:
\begin{lstlisting}
inductive Amt where
  | const (n : Nat)
  | per (w : Whose) (k : Nat)
  | add (a b : Amt)
  | sub (a b : Amt)
  | many (a : Amt)
\end{lstlisting}
\noindent \lean{per w k} is $k$ for each damage counter on the \lean{w} Active \pkmn{},
\lean{sub} stops at 0, and \lean{many a} is the sum of any number of values of \lean{a}, for
``flip a coin until you get tails; 30 damage for each heads''. Jynx's ``20 damage plus 10 more
for each damage counter on the Defending \pkmn{}'' (\texttt{base1}-31) is
\lean{add (const 2) (per opp 1)}. An amount is evaluated once, when the damage step is reached
(\lean{Amt.Val}, a relation because of \lean{many}). The effects are:

\begin{lstlisting}
inductive Effect where
  | skip | draw
  | heal (n : Nat)        -- remove up to n counters from the player's Active
  | counters (n : Nat)    -- place n counters on the opponent's Active
  | inflict (c : Condition) -- apply a Special Condition to the opponent's Active
  | poisonUp (n : Nat)    -- if the opponent's Active is Poisoned, n more Poison counters
                          -- at every Checkup; repeated uses add up
  | cure                  -- the player's Active recovers from all Special Conditions
  | shuffle | search      -- shuffle own deck; move one chosen card from deck to hand
  | handToBottom          -- put the hand on the bottom of the deck in a chosen order
  | ite (t : Test) (a b : Effect)  -- if the test holds then a else b
  | scale (w : Whose) (k : Nat)    -- k counters on the opponent's Active for each
                                   -- damage counter on the w Active
  | selfDamage (n : Nat)  -- place n counters on the player's own Active
  | healAll               -- remove all damage counters from the player's own Active
  | mark (w : Whose) (m : Marker)   -- put marker m on one of the w player's Pokémon in play
  | unmark (w : Whose) (m : Marker) -- remove every m marker from all of the w player's Pokémon
  | color (w : Whose) (t : EType)   -- Coloring counter for type t on one of the w player's
                                    -- Pokémon in play, replacing any earlier one
  | linger (w : Whose) (s : Slot) (m : Mod) -- the w player gets the lingering modifier m
                                            -- for the turns s
  | damage (a : Amt)      -- the attack's damage step: base damage a, through the damage
                          -- calculation, to the opponent's Active
  | drain (a : Amt)       -- as damage a, and heal the player's Active by the damage done
  | seq (a b : Effect)
  | choice (a b : Effect)
  | iter (a : Effect)
\end{lstlisting}
(The comments abbreviate the docstrings, and \lean{skip}, \lean{draw}, \lean{shuffle} and
\lean{search} are written two to a line; constructors and types are as in \texttt{Card.lean}.)
\lean{Effect} has twenty-three constructors. \texttt{Effect.lean} defines
\lean{Effect.isPrim}, which is false exactly for \lean{ite}, \lean{seq}, \lean{choice} and
\lean{iter}. We call these four the \emph{combinators} and the other nineteen
(\lean{skip} included) the \emph{primitives}; a primitive resolves in one step with no
sub-derivation. \lean{choice} covers both coin flips and player decisions, and \lean{iter}
covers ``flip until tails'' and ``as often as you like''. The primitive \lean{poisonUp} models
Croagunk's Poison Up (``This effect can be applied more than once''). \lean{scale opp 1}
doubles the damage counters on the opponent's Active \pkmn{}, as Lunala (\texttt{cel25}-15)
does, and \lean{scale} also covers attacks that place damage counters ``for each damage counter
on'' an Active \pkmn{}.

The relation \lean{Run e g g'} holds when resolving \lean{e} in state \lean{g} can produce
\lean{g'}. It is nondeterministic: choices, coin flips and shuffles are all branches. This
suits reachability questions and is silent about probability. \lean{Run} is the semantics of
an attack's effect, and the rulebook checks for Knock Outs only at the end of an attack, so
nothing is Knocked Out inside one \lean{Run}: damage may exceed HP, and the tests and
\lean{scale} see the excess; Appendix~\ref{app:defs} shows the cases \lean{iteT} and
\lean{scale}. A theorem about the model speaks only about decks
whose card text is expressible in \lean{Effect}. Tests and \lean{scale} read the state only
through the two Active \pkmn{}: their damage, their Special Conditions, their markers and their
type. The damage step (\lean{damage} or \lean{drain}) reads the same two \pkmn{}: its amount
may read the damage on either one, and its calculation reads the attacker's type, the
defender's Weakness and Resistance, and both players' lingering modifiers in effect
(Section~\ref{sec:damage-calc}). No effect reads Poison strength, HP, the Bench,
Energy, or the cards in any hand, deck, discard pile or Prize list, and no test reads a
modifier directly. The marker effects are the only ones that write to a Benched \pkmn{}.

\subsection{Special Conditions and Checkup}

Asleep and Paralyzed \pkmn{} can neither attack nor retreat. A Confused \pkmn{} may try to
attack: on heads the attack happens (\lean{attack}), and on tails it does not, and the attacker
takes 3 damage counters (\lean{attackConfused}); the failed attempt still ends the turn. The
rulebook says only that ``the attack doesn't happen'' \cite[p.~16]{rb-151}; the model reads the
Confusion check as step C of using the attack \cite[p.~20]{rb-151}. Retreating and evolving
remove every Special Condition and reset the Poison strength to 0, and so does a new Poison,
which replaces the old one (\lean{InPlay.inflict}). That Croagunk's raised Poison strength ends
with the replaced Poison is the model's reading of ``replaces''; we found no ruling on it.
\pkmn{} Checkup applies to one Active \pkmn{} as
follows (in namespace \lean{InPlay}); the last term is the Char marker of
Section~\ref{sec:markers}:

\begin{lstlisting}
def checkup (ownersTurnEnded burnHeads sleepHeads charTails : Bool) (p : InPlay) : InPlay :=
  { p with
    damage := p.damage + (if p.poisoned then 1 + p.poisonUp else 0) + (if p.burned then 2 else 0)
      + (if p.marks.char && charTails then 2 else 0),
    burned := p.burned && !burnHeads,
    rot := match p.rot with
      | .asleep => if sleepHeads then .none else .asleep
      | .paralyzed => if ownersTurnEnded then .none else .paralyzed
      | r => r }
\end{lstlisting}

The move \lean{endTurn} applies Checkup to both Active \pkmn{}, swaps the players, ages both
players' lingering modifiers (\lean{Store.shift}, Section~\ref{sec:linger-mech}), and enters
a \emph{between-turns} phase (\lean{between = true}); Appendix~\ref{app:defs} gives its
definition. Between turns only Knock Outs, promotions, and the draw that begins the \emph{turn
proper} (\lean{startTurn}, or \lean{deckOut} if the deck is empty) are legal; the turn proper
is the rest of the turn, after that draw. No effect resolves between turns.

\subsection{Markers}\label{sec:markers}

A \emph{marker} is a physical token that a few old cards place on a \pkmn{}. The model keeps
four markers as flags (Imprison, Shock-wave, Lightning Rod and Char) and the Coloring marker,
which changes a \pkmn{}'s type, as an optional type:
\begin{lstlisting}
structure Marks where
  imprison : Bool
  shockwave : Bool
  lightningRod : Bool
  char : Bool
  /-- the type set by a Coloring counter, if the Pokémon has one -/
  coloring : Option EType
\end{lstlisting}
The effects \lean{mark}, \lean{unmark} and \lean{color} write these fields, the tests
\lean{marked} and \lean{isType} read them on either Active \pkmn{}, and Checkup puts 2 damage
counters on an Active \pkmn{} with a Char marker when the flip is tails.

The printed Imprison and Shock-wave markers have no cap: a card places one at a time, and
nothing limits how many a \pkmn{} carries. Our survey read every sentence that mentions a
marker (Section~\ref{sec:findings}). The only reads are ``has any Imprison markers''
(Gardevoir ex $\delta$) and ``has any Shock-wave markers'' (Tyranitar ex $\delta$), and the only
removal takes them all (Tropius $\delta$). Lightning Rod has a printed cap of one, and a Char
marker is placed only on a \pkmn{} without one. \texttt{Markers.lean} records why a flag is
enough:
\begin{lemma}[\lean{Markers.flag\_exact}]\label{thm:flag}
For every sequence of placements and removals, ``has any'' gives the same answer on the count
as on the flag. (The Lean statement is in Appendix~\ref{app:defs}.)
\end{lemma}
\noindent The lemma concerns an abstract count with these two operations and this one read; it
is not a statement about \lean{Move}. That the printed cards perform only these operations is
the survey's finding, for English card text only, and a stated assumption of the model: a card
that counted these markers, or removed one at a time, would break it. Markers are finitely many
flags per \pkmn{}, they are not clamped, and the finite list of Theorem~\ref{thm:finite}
enumerates them. \lean{Move.evolve\_marks} and \lean{Move.retreat\_marks} show that evolving
and retreating keep them. The departures that concern markers are item~10 of
Appendix~\ref{app:departures}.

\subsection{Lingering modifiers}\label{sec:linger-mech}

Each player has a \emph{store} of lingering modifiers, kept by \emph{slot}, the turn in which
they are in effect: this turn (\lean{now}), the next turn, which is the opponent's
(\lean{next}), the turn after that (\lean{after}), and the rest of the game (\lean{always}).
We call the \lean{always} slot \emph{rest-of-game}, and modifiers in the \lean{next} and
\lean{after} slots, which are written in one turn and in effect in a later one,
\emph{carried}. A modifier is a damage bonus for the player's attacks, a damage reduction or a
prevention of damage for attacks against the player's \pkmn{}, or a lock on attacking,
retreating or playing Items. In a store, amounts add and flags combine:
\begin{lstlisting}
inductive Slot where
  | now | next | after | always

structure Mods where
  bonus : Nat
  reduce : Nat
  prevent : Bool
  noAttack : Bool
  noRetreat : Bool
  noItem : Bool

structure Store where
  now : Mods
  next : Mods
  after : Mods
  always : Mods

def live (t : Store) : Mods := t.now.merge t.always
def shift (t : Store) : Store := ⟨t.next, t.after, Mods.empty, t.always⟩
\end{lstlisting}
\noindent (\lean{live} and \lean{shift} are in namespace \lean{Store}; \lean{merge} adds the
amounts and takes the disjunction of the flags.) The move \lean{endTurn} applies
\lean{Store.shift} to both players' stores, so this turn's modifiers expire and the others
move one turn closer. Four lemmas state the result; the last, \lean{Store.add\_always\_shift}
(Appendix~\ref{app:defs}), says that a rest-of-game modifier survives a shift. Three
primitives use the stores. \lean{linger w s m} gives the \lean{w} player the modifier \lean{m}
in slot \lean{s}. \lean{damage a} does damage to the opponent's Active \pkmn{} through the
damage calculation of Section~\ref{sec:damage-calc}, which reads the modifiers in effect, and
\lean{drain a} does the same and heals the player's Active \pkmn{} by the damage done. Locks
add a premise to \lean{attack}, \lean{attackConfused}, \lean{retreat} and \lean{item}.

\lean{linger} is a general primitive, usable any number of times and inside \lean{iter}, so
bonuses and reductions accumulate without limit, and a rest-of-game modifier, which printed
cards grant only through once-per-game attacks, can be granted by an ordinary attack used
every turn, and to either player. We chose this over-approximation so that the model does not
rest on the empirical bound on printed modifiers of Section~\ref{sec:linger-printed}; the
faithful theorem of Section~\ref{sec:faithful} recovers a bound from a condition on the deck
instead. The universal construction of Section~\ref{sec:linger} uses the over-approximation.

\subsection{Costs and the damage calculation}\label{sec:damage-calc}

An attack's cost is a list of types. Each attached Energy card provides one unit of its type
(\lean{Card.provides}, over \lean{InPlay.units}), and the cost is paid when its symbols can be
matched to distinct units, a typed symbol to a unit of that type and a Colorless symbol to any
unit (\lean{EType.paidBy}):
\begin{lstlisting}
inductive Pays : List EType → List (List EType) → Prop where
  | nil (us : List (List EType)) : Pays [] us
  | cons {s : EType} {cs : List EType} {u : List EType} {us₁ us₂ : List (List EType)}
      (hu : s.paidBy u = true) (h : Pays cs (us₁ ++ us₂)) : Pays (s :: cs) (us₁ ++ u :: us₂)
\end{lstlisting}
\noindent This is a matching, not a count: two Fire symbols are not paid by one Fire and one
Water Energy (\lean{pays\_fire\_fire}). The function \lean{payable} decides \lean{Pays}
(\lean{payable\_iff}). A Retreat Cost is paid by discarding that many Energy cards, which with
Basic Energy only is the rulebook's rule.

The damage calculation follows the rulebook's steps \cite[p.~20]{rb-151}, for an attack whose
base damage is \lean{n} counters:
\begin{lstlisting}
def dealtBy (n : Nat) (weak : Bool) (resist : Nat) (own opp : Store) : Nat :=
  if n = 0 then 0 else if opp.live.prevent then 0 else
    ((if weak then 2 * (n + own.live.bonus) else n + own.live.bonus) - resist) - opp.live.reduce
\end{lstlisting}
\noindent with truncated subtraction. \lean{GameState.dealt} applies it with the Weakness and
Resistance of the opponent's Active \pkmn{} to the type of the player's (a Coloring marker
changes that type) and with both players' stores. An attack with base damage 0 does no damage,
whatever the modifiers (\lean{dealtBy\_zero}). Weakness doubles the base damage together with
the bonus, since the bonus is added first; the rulebook's example, Fuecoco's 10 damage against
a Sprigatito weak to Fire, gives 2 damage counters (\lean{Calc.dealtBy\_gnaw}). A stored bonus
enters before Weakness and Resistance and a stored reduction after them, which is what printed
modifiers say (``more damage \dots\ (before applying Weakness and Resistance)'', ``less damage
\dots\ (after applying Weakness and Resistance)''); a modifier that names the other step is not
expressible.

An attack makes the calculation at most once. What it resolves, and when it is well formed, are
\begin{lstlisting}
def Attack.script (atk : Attack) : Effect :=
  if atk.damage = 0 then atk.effect else .seq (.damage (.const atk.damage)) atk.effect

def Attack.wf (atk : Attack) : Bool :=
  atk.effect.oneHit && (decide (atk.damage = 0) || atk.effect.hitFree)
\end{lstlisting}
\noindent An attack with a printed base damage and no damage text does that damage and then its
effect. An attack whose text computes or conditions its damage has base damage field 0 and one
damage step in its effect; effects before that step are the rulebook's ``effects that happen
before damage'', and effects after it see the damage placed. \lean{Effect.hitFree} says that an
effect has no damage step, and \lean{Effect.oneHit} that it has at most one on each line of
resolution: one in each branch of \lean{ite} and \lean{choice}, one in at most one part of a
\lean{seq}, and none under \lean{iter}. A damage that depends on a test, a coin flip or a
choice is written with the damage step inside the branches. \lean{Move.attack} requires
\lean{Attack.wf}, and \lean{Move.item} and \lean{Move.supporter} require a \lean{hitFree}
effect: printed Trainer cards place damage counters, not attack damage.

This design gives up one generality: an attack cannot make several damage steps, such as an
\lean{iter} of \lean{damage} that applies the modifiers at every step, which is not the
rulebook's calculation, and an Item or Supporter cannot do attack damage. Two damage
calculations in one attack (an attack that also damages a Benched \pkmn{}) and damage counted
from anything other than damage counters on an Active \pkmn{} (Energy, the Bench, cards in hand
or in the discard pile) are not expressible. What it gains is the rulebook's calculation on the
whole base damage: with 3 damage counters on the target and a reduction of 3, Jynx's attack
does $(2 + 3) - 3 = 2$ (\lean{Calc.dealtBy\_jynx}), whereas placing the scaled part as damage
counters outside the calculation would give 3.

\subsection{Knock Out timing and moves}\label{sec:moves}

The rulebook checks for Knock Outs at the end of every attack and at the end of every Checkup,
and the Advanced Player's Rulebook also at once whenever a \pkmn{}'s damage reaches its HP at any
other time (Section~\ref{sec:game}). The model follows both with two definitions.

\begin{definition}[\lean{DamageBounded}, \lean{Settled}]
No Knock Out is \emph{pending} if every \pkmn{} in play has fewer damage counters than its HP.
A state is \emph{settled} if no Knock Out is pending and both Active Spots are filled.
\begin{lstlisting}
def DamageBounded (g : GameState) : Prop :=
  ∀ p : InPlay,
    (g.cur.active = some p ∨ p ∈ g.cur.bench ∨ g.opp.active = some p ∨ p ∈ g.opp.bench) →
      p.damage < p.top.hp

def Settled (g : GameState) : Prop :=
  DamageBounded g ∧ g.cur.active ≠ none ∧ g.opp.active ≠ none
\end{lstlisting}
\end{definition}

\lean{Move g g'} is one legal step for the player to move. Six moves resolve a Knock Out and
are legal in either phase: \lean{knockOut} and \lean{knockOutOwn} (an Active \pkmn{} at or
above its HP goes to the discard pile, and the other player takes the top Prize card, winning
if it was the last), \lean{promote} and \lean{promoteOwn} (a Benched \pkmn{} fills an empty
Active Spot), and \lean{noPokemon} and \lean{noPokemonOwn} (the side with no \pkmn{} left
loses). Every other move requires the state to be settled (\lean{endTurn}, \lean{startTurn}
and \lean{deckOut} state the same condition as \lean{DamageBounded} with both Active \pkmn{}
present). So once a move leaves a \pkmn{} at or above its HP, the next moves must Knock it Out
and promote a replacement before anything else happens, and no move reads, heals or otherwise
uses a \pkmn{} above its HP. During the turn proper the other moves are \lean{playBasic} (a
Basic \pkmn{} with positive HP to the Bench, Bench below 5), \lean{attach} (one Energy per turn
to the Active), \lean{evolve} (the Active, not on the player's own first turn, that is while
\lean{hadTurn} is false, and not if put into play or evolved this turn), \lean{retreat} (pay
the cost in Energy, swap with a Benched \pkmn{}, not Asleep or Paralyzed), \lean{item},
\lean{supporter} (once per turn, not on the first turn of the game), \lean{attack} (not on the
first turn of the game; the cost paid by \lean{Pays}; the attack well formed; its script
resolved by \lean{Run}), \lean{attackConfused}, and \lean{endTurn}, which sets \lean{hadTurn}
for the player whose turn ends. An attack, or a failed Confused attempt, sets
\lean{attacked}, and every action of the turn proper requires \lean{attacked} to be false, so
after an attack only its Knock Outs and \lean{endTurn} are legal. An attack carries a flag
\lean{gx}; a GX attack, or a Confused attempt at one, needs \lean{gxUsed} to be false, and a GX
attack that is made sets it, so each player makes at most one GX attack per game. A failed
Confused attempt does not set it; the model takes this answer to an unverified ruling. Between
turns the other moves are \lean{startTurn} (draw) and \lean{deckOut} (the deck is empty; the
player to move loses).

\paragraph{The turn.} The turn flags of a start state are \emph{consistent} if, on the first
turn of the game, the player to move went first, neither player has had a turn and nobody has
attacked, and, between turns, nobody has attacked. Every start state of a game is consistent;
the model does not require it of a start state, but play preserves it.
\begin{lstlisting}
def GameState.FlagsOk (g : GameState) : Prop :=
  (g.firstTurn = true →
    g.curIsFirst = true ∧ g.cur.hadTurn = false ∧ g.opp.hadTurn = false ∧ g.attacked = false) ∧
  (g.between = true → g.attacked = false)
\end{lstlisting}

\begin{theorem}[\lean{Move.after\_attack}, \lean{Reach.flagsOk} and two
companions]\label{thm:turn}
An attack ends the turn: once the player to move has attacked, every move is a Knock Out move or
passes the turn. Play preserves consistent flags. With consistent flags, no move on the first
turn of the game leaves the player to move having attacked, and neither player has had a turn,
so \lean{evolve} is not legal.
\begin{lstlisting}
theorem Move.after_attack {g g' : GameState} (hm : Move g g') (ha : g.attacked = true)
    (hb : g.between = false) :
    KOMove g g' ∨ (g'.curIsFirst = !g.curIsFirst ∧ g'.between = true ∧ g'.attacked = false)
theorem Reach.flagsOk {g₀ g : GameState} (h : Reach g₀ g) (hc : g₀.FlagsOk) : g.FlagsOk
theorem Move.firstTurn_no_attack {g g' : GameState} (hc : g.FlagsOk) (hf : g.firstTurn = true)
    (hm : Move g g') : g'.attacked = false
theorem GameState.FlagsOk.firstTurn_no_evolve {g : GameState} (hc : g.FlagsOk)
    (hf : g.firstTurn = true) : g.cur.hadTurn = false ∧ g.opp.hadTurn = false
\end{lstlisting}
\end{theorem}
\noindent \lean{Move.after\_attack\_passTurn} sharpens the first statement: a move after an
attack that is not a Knock Out move gives exactly the state that ending the turn gives, for some
Checkup coin flips. \lean{Move.firstTurn\_no\_supporter} states the Supporter ban in the same
way. The evolution theorem covers the first turn of the game only; the ban on the second
player's own first turn is the \lean{hadTurn} premise of \lean{evolve} itself, not a theorem
under \lean{FlagsOk}, which constrains \lean{hadTurn} only while \lean{firstTurn} holds.

\paragraph{Item and Supporter effects.} We call the effect of an Item or Supporter card a
\emph{Trainer effect}; Tools and Stadiums have none in the model. A Trainer effect has no damage
step, and it is resolved not by \lean{Run} but by \lean{TRun}, which takes every Knock Out at
once after each primitive step, following the Advanced Player's Rulebook
(Section~\ref{sec:game}).
\lean{KOMove} is the relation of the six Knock Out moves, with the same premises and results
(\lean{KOMove.move}), and \lean{Resolve} applies them until the state is settled or the game
has a winner. A primitive step of \lean{TRun} is a \lean{Run} of the primitive followed by
\lean{Resolve}; the combinators have the same shape as in \lean{Run}, and no step is taken once
the game has a winner (Appendix~\ref{app:defs}). So a Knock Out caused by one step gives its Prize card and
promotes the replacement, chosen nondeterministically, before the next step; the rest of the
effect acts on the new board; and once the game has a winner the rest of the effect is not
resolved. The moves \lean{item} and \lean{supporter} start \lean{TRun} from a settled state.

Reachability is the reflexive-transitive closure of \lean{Move}, stopped once a winner is set.
A move is made by whichever player is to move, so \lean{Reach g$_0$ g} says that \emph{some}
sequence of legal moves of the two players leads from \lean{g$_0$} to \lean{g}:

\begin{lstlisting}
inductive Reach : GameState → GameState → Prop where
  | refl (g : GameState) : Reach g g
  | step {g₀ g g' : GameState} : Reach g₀ g → g.winner = none → Move g g' → Reach g₀ g'
\end{lstlisting}

\begin{definition}[\lean{Benched}, \lean{WellFormed}]
A Benched \pkmn{} is below its HP and has no Special Condition. A state is \emph{well formed}
if, while there is no winner, both players have Prize cards left, and every Benched \pkmn{} is
\lean{Benched}.
\begin{lstlisting}
def Benched (p : InPlay) : Prop :=
  p.damage < p.top.hp ∧ p.rot = .none ∧ p.poisoned = false ∧ p.burned = false ∧ p.poisonUp = 0

def WellFormed (g : GameState) : Prop :=
  (g.winner = none → g.cur.prizes ≠ [] ∧ g.opp.prizes ≠ []) ∧
    (∀ p ∈ g.cur.bench, Benched p) ∧ ∀ p ∈ g.opp.bench, Benched p
\end{lstlisting}
\end{definition}

\subsection{Tame decks}\label{sec:tame}

The first finiteness theorem needs a hypothesis on the deck. An effect is \emph{tame} if it adds no damage
reduction and uses no \lean{drain}; a card is tame if its effects are. The hypothesis
\lean{Fit k} on a state collects what the clamp with cap $k$ (Section~\ref{sec:clamp}) needs:
\begin{lstlisting}
def Effect.tame : Effect → Bool
  | .linger _ _ (.reduce _) => false
  | .drain _ => false
  | .ite _ a b => a.tame && b.tame
  | .seq a b => a.tame && b.tame
  | .choice a b => a.tame && b.tame
  | .iter a => a.tame
  | _ => true

def GameState.Fit (k : Nat) (g : GameState) : Prop :=
  0 < k ∧ g.HpLe k ∧ (∀ x ∈ g.allCards, x.tame = true ∧ x.tmax < k) ∧ g.NoReduce
\end{lstlisting}
\noindent Here \lean{HpLe k} says that every printed HP and every printed Resistance among the
cards is at most $k$,
\lean{x.tmax} is the largest constant with which the card's tests compare damage, and
\lean{NoReduce} says that no damage reduction is in effect in either player's store.
\lean{Fit} is preserved by play (\lean{Reach.fit}).

\subsection{Departures from the rulebook}\label{sec:departures}

The header of \texttt{Rules.lean} lists every known place where the model is simpler than the
rulebook, and the sources of the turn and damage rules. Appendix~\ref{app:departures} restates
both, the list in full and in its order as items 1--23; Table~\ref{tab:departures} groups it. We classify each group by its effect on the
finiteness theorems; the classification is an argument, not mechanized. A departure is
\emph{load-bearing} if it makes finiteness easier to prove than it would be for the rulebook,
so that the theorem does not transfer to the printed game until the departure is removed and
the theorem re-proved. It is \emph{neutral or harder} if it gives the model more behaviours
than the printed game, or replaces one finite choice by another. A departure that
\emph{restricts expression} removes card texts from the language, and a theorem about the model
says nothing about decks that need them. We do not claim that any unclassified departure is
neutral.

\begin{table}[ht]
\centering
\footnotesize
\begin{tabularx}{\textwidth}{@{}>{\raggedright\arraybackslash}X l >{\raggedright\arraybackslash}p{4.3cm}@{}}
\toprule
Departure & Items & Effect on finiteness \\
\midrule
No set-up; start states need not be well formed or have consistent turn flags & 1, 22 & not classified \\
Energy and evolution on the Active only & 2 & neutral or harder \\
Basic Energy only; no Special Energy or cost changes & 3 & restricts expression \\
Weakness $\times 2$ and Resistance $-N$ to one type only; no effects on Weakness or Resistance & 4 & restricts expression \\
Knock Outs as separate moves, one at a time; Active only; one Prize, the top one & 5, 8 & not classified \\
Knock Outs at other times rest on the Advanced Player's Rulebook; in Item and Supporter effects after each primitive, with immediate promotion & 6, 7 & argued not to make finiteness easier \\
HP is the printed value & 9 & load-bearing \\
Markers as flags, under an assumption about the card pool & 10 & load-bearing (the assumption) \\
Effects read only the two Active \pkmn{}; no test on Poison strength or HP & 11 & load-bearing (the Poison test); otherwise restricts expression \\
Lingering modifiers stack without limit and belong to a player; no ``attack twice'' or extra turns & 12 & load-bearing; met by the tameness or faithfulness hypothesis, and the absences by the faithful bound \\
One damage step per attack, over the amount language; no attack damage in Trainer effects & 13 & restricts expression \\
General primitives (\lean{ite}, \lean{scale}, \lean{selfDamage}, \lean{healAll}, \lean{poisonUp}) & 14, 18 & neutral or harder \\
GX attacks once per game; a failed Confused GX attempt does not use it up & 15 & not classified \\
Special Conditions and Checkup simplified; coin flips nondeterministic & 16, 17 & item~17 neutral or harder; item~16 not classified \\
Damage and healing only on the two Active \pkmn{}; no Abilities; Stadiums and Tools inert & 19--21 & not classified; restricts expression \\
Card text that lifts a first-turn ban is not expressible & 23 & restricts expression \\
\bottomrule
\end{tabularx}
\caption{The model's departures from the rulebook, grouped; Appendix~\ref{app:departures}
gives each item.}\label{tab:departures}
\end{table}

\paragraph{Load-bearing departures.} Four departures are load-bearing. HP is the printed value,
so the cap and the finite list are built from printed HP. Markers are flags, and what is
load-bearing is the assumption about the card pool behind Lemma~\ref{thm:flag}. Lingering
modifiers stack without limit, and each finiteness theorem meets this by a hypothesis on the
deck, tameness or faithfulness (Section~\ref{sec:finite-side}); the faithful bound also rests on
three absences in the model, of card recycling, of repeated attacks and of extra turns
(Section~\ref{sec:faithful}). And no effect tests Poison strength: both clamps cap Poison
strength at HP, and a test against a larger number would observe the excess; no part of our
survey searched for such a test.

\paragraph{Knock Out timing inside Item and Supporter effects.} No English rulebook says when
Knock Outs are checked outside attacks and Checkup, so the model's timing there rests on the
Advanced Player's Rulebook of \pkmn{} Asia; it is a modelling decision until an English source
confirms it. Checking for Knock Outs after each primitive step makes the model's timing that
document's, up to two caveats, each of which can lose printed outcomes. First, a printed sentence that is atomic but that the model writes
as several primitives gets a check between them, so the outcome in which the whole sentence
acts on the original \pkmn{} is absent: \lean{seq (counters 5) (scale opp 1)} against a 5-HP
Active always acts on the promoted \pkmn{} in its second step. Second, the replacement is
promoted at once, so outcomes in which later steps find an empty Active Spot are absent; we
found no ruling on when the replacement is chosen in mid-effect. We argue that neither caveat
makes finiteness easier, since the atomic reading is the attack case, which \lean{Run} already
covers, and deferred promotion is a finite-branching variant. The argument is not mechanized,
and no rulebook is formalized, so the agreement is not a theorem.

%======================================================================
\section{The finite side}\label{sec:finite-side}
%======================================================================

This section proves two finiteness theorems: the clamped game of a tame deck, and the
exact-store clamped game of a faithful deck, are finite and have the same moves as the game
itself. It first shows that the cards are conserved and that, without a clamp, the reachable
states are infinite.

\subsection{Card conservation}

\begin{theorem}[\lean{Run.perm}, \lean{TRun.perm}, \lean{Move.perm}, \lean{Reach.perm}]\label{thm:perm}
Resolving an effect, resolving a Trainer effect with its Knock Outs, making a legal move, and
legal play in general permute the cards in the game. (Lean statements:
Appendix~\ref{app:defs}.)
\end{theorem}
The corollary \lean{Reach.cards\_eq} states that the number of cards is constant. Counting
alone would not suffice for finiteness, because \lean{Card} is an infinite type (names are
strings, HP is a natural number); the permutation statement confines every reachable state to
rearrangements of the starting cards.

\subsection{Without a clamp the state space is infinite}

Two numeric components can be made large by legal play.

\begin{theorem}[\lean{counters\_unbounded}, \lean{poisonUp\_unbounded}]\label{thm:unbounded}
Inside one effect, the damage on the opponent's Active \pkmn{}, and the Poison strength of a
Poisoned opposing Active \pkmn{}, can reach any value: for every $n$, the effects
\lean{iter (counters 1)} and \lean{iter (poisonUp 1)} have resolutions that add exactly $n$
(Appendix~\ref{app:defs}).
\end{theorem}

The effect \lean{iter (counters 1)} models ``flip a coin until you get tails; one damage
counter for each heads''. The excess damage is removed at the next Knock Out, but Poison
strength is not: Poison does no damage until Checkup, so a large Poison strength can sit on a
\pkmn{} that is below its HP. \texttt{Infinite.lean} turns this into a statement about
reachable states. Its start state \lean{Infinite.start} gives each player one Prize card and a
50-HP Active \pkmn{}, and gives the player to move one Item whose effect is

\begin{lstlisting}
def eff : Effect := .seq (.inflict .poisoned) (.iter (.poisonUp 1))
\end{lstlisting}
\noindent Playing it reaches, for every $n$, a state \lean{st n} with Poison strength $n$ on
the opponent's Active and no pending Knock Out.

\begin{theorem}[\lean{Infinite.not\_finite} and two companions]\label{thm:infinite}
The unclamped reachable state space is infinite, and so is its subset with no pending Knock
Out, from a start that is well formed and has no pending Knock Out.
\begin{lstlisting}
theorem not_finite : ¬ ∃ L : List GameState, ∀ g, Reach start g → g ∈ L
theorem not_finite_damageBounded :
    ¬ ∃ L : List GameState, ∀ g, Reach start g → DamageBounded g → g ∈ L
theorem _root_.Pokemon.Reach.infinite_damageBounded :
    ∃ g₀, WellFormed g₀ ∧ DamageBounded g₀ ∧
      ¬ ∃ L : List GameState, ∀ g, Reach g₀ g → DamageBounded g → g ∈ L
\end{lstlisting}
(The first two are in namespace \lean{Pokemon.Infinite}, where \lean{start} is the state
described above; \lean{Infinite.start\_wellFormed} and \lean{Infinite.start\_damageBounded}
show that it is well formed and has no pending Knock Out.)
\end{theorem}

Theorem~\ref{thm:infinite} rests on a generality of the model: the printed Croagunk raises
Poison strength once per use of its attack, not without bound inside one effect
(Appendix~\ref{app:departures}, item~18). Theorem~\ref{thm:unbounded} does not: the rulebook
itself allows a transient excess of damage within an attack (Section~\ref{sec:game}). In both
cases the reachable states are unbounded, and a finiteness theorem needs a clamp.

\begin{remark}[pending Knock Outs between turns]\label{rem:pendingko}
Checkup can leave a pending Knock Out when the turn passes, in play reached from a well-formed
start (\lean{Infinite.endTurn\_pending\_ko\_reachable}; the witness is \lean{st 5}, on which
Checkup places $1 + 5$ Poison damage counters on a 50-HP Active). The Knock Out is taken
between turns, since \lean{startTurn} requires that none be pending:
\begin{lstlisting}
theorem Move.damageBounded_of_startTurn {g g' : GameState} (h : Move g g')
    (hb : g.between = true) (hb' : g'.between = false) : DamageBounded g'
\end{lstlisting}
This matches the rulebook, which checks for Knock Outs at the end of Checkup.
\end{remark}

\subsection{Progress and Knock Out timing}\label{sec:kotiming}

\begin{theorem}[\lean{Reach.wellFormed}, \lean{Reach.progress},
\lean{Reach.knockOut\_of\_pending}]\label{thm:progress}
Legal play from a well-formed start stays well formed; no reachable state without a winner is
stuck; and if a Knock Out is pending, some move Knocks Out an Active \pkmn{} and takes one
Prize card. (Lean statements: Appendix~\ref{app:defs}.)
\end{theorem}
\lean{Reach.progress} says only that some move exists; since \lean{endTurn} is available
whenever nothing is pending, it does not say that the game ends.

Between moves, the Knock Out timing follows from \lean{Settled} being a premise of every move
other than the six that resolve a Knock Out.

\begin{theorem}[\lean{Move.pending\_untouched}]\label{thm:untouched}
A move leaves an Active \pkmn{} at or above its HP either unchanged in its Active Spot or
removed from it (Knocked Out), and does not pass the turn. (Lean statement:
Appendix~\ref{app:defs}.)
\end{theorem}

Inside an Item or Supporter effect, the timing is stated by comparing \lean{TRun} with a
relation \lean{TRunG}, defined only for this purpose, which is \lean{TRun} with the extra
premise \lean{Settled g} on every primitive step and every test (on the constructors
\lean{prim}, \lean{iteT} and \lean{iteF}).

\begin{theorem}[\lean{TRun.iff\_guarded}, \lean{TRun.settled\_or\_won}]\label{thm:guarded}
From a state that is settled or finished, as every \lean{item} and \lean{supporter} move
starts, \lean{TRun} and \lean{TRunG} coincide: no primitive step of a Trainer effect is
executed, and no test is evaluated, while a \pkmn{} in play is at or above its HP or an Active
Spot is empty. Such an effect ends settled or with the game over. (Lean statements:
Appendix~\ref{app:defs}.)
\end{theorem}
\noindent \texttt{Trainer.lean} also proves three witnesses, each a single legal move from a
well-formed, settled start, of behaviours that resolving an Item as a whole would not produce:
a Knock Out of the player's own Active in mid-effect, a Knock Out followed by damage to the
promoted \pkmn{}, and a win in mid-effect after which the rest of the effect is not resolved
(Theorem~\ref{thm:witness} in Appendix~\ref{app:audit}).

\subsection{The clamp}\label{sec:clamp}

\begin{definition}[\lean{Cap}, \lean{InPlay.clampWith}, \lean{Store.cap},
\lean{GameState.clampAt}, \lean{GameState.clamp}, \lean{GameState.rclamp}]
\begin{lstlisting}
inductive Cap where
  | exact
  | hp
  | upto (k : Nat)

def Cap.dmg : Cap → InPlay → Nat
  | .exact, p => p.damage
  | .hp, p => min p.damage p.top.hp
  | .upto k, p => min p.damage k

def InPlay.clampWith (c : Cap) (p : InPlay) : InPlay :=
  { p with damage := c.dmg p, poisonUp := min p.poisonUp p.top.hp }

def Mods.cap (k : Nat) (m : Mods) : Mods :=
  { m with bonus := min m.bonus k, reduce := min m.reduce k }

def Store.cap (k : Nat) (t : Store) : Store :=
  ⟨t.now.cap k, t.next.cap k, t.after.cap k, t.always.cap k⟩

def GameState.clampAt (co cp : Cap) (k : Nat) (g : GameState) : GameState :=
  { g with cur := { g.cur with active := g.cur.active.map (InPlay.clampWith co),
                               mods := g.cur.mods.cap k },
           opp := { g.opp with active := g.opp.active.map (InPlay.clampWith cp),
                               mods := g.opp.mods.cap k } }

def GameState.clamp (k : Nat) (g : GameState) : GameState := g.clampAt .hp .hp (2 * k)

def GameState.rclamp (k : Nat) (g : GameState) : GameState := g.clampAt .exact (.upto k) (2 * k)
\end{lstlisting}
\end{definition}

The clamp \lean{clamp k} caps the damage and the Poison strength of each Active \pkmn{} at its
printed HP, in every phase, and caps every damage bonus and reduction in both players' stores
at $2k$. Every other component is unchanged, including the damage on every Benched \pkmn{} and
the lingering flags. Two states \emph{have the same clamp} if their images under
\lean{clamp k} are equal. \lean{rclamp k} is the map used inside an attack's effect: the
player's own damage exact, the opponent's capped at $k$, the stores at $2k$. The cap that
matters is
\begin{lstlisting}
def bound (C : List Card) : Nat := maxHp C + maxResist C + cardsTmax C + 1
\end{lstlisting}
\noindent which exceeds every printed HP, every printed Resistance and every test constant among
the cards $C$. Table~\ref{tab:clamp} gives the reason no move can observe each excess, under the
hypothesis \lean{Fit k} of Section~\ref{sec:tame}.

\begin{table}[ht]
\centering
\small
\begin{tabularx}{\textwidth}{@{}l l >{\raggedright\arraybackslash}X@{}}
\toprule
Quantity & Capped at & Why no move observes the excess \\
\midrule
Damage, either Active & HP & Every move other than a Knock Out move starts settled, where damage is below HP and the clamp changes nothing. A \pkmn{} at or above its HP is read only by the Knock Out test \lean{hp $\le$ damage} and its negation, which the clamp preserves. \\
Poison strength, either Active & HP & Read only by Checkup, which adds $1 + {}$\lean{poisonUp} to damage; $\min(d + 1 + u, h)$ depends on $u$ only through $\min(u, h)$. No test reads it. Effects only add to it; curing, retreating and evolving reset it to 0. \\
Damage bonus, both stores & $2k$ & Read only by the damage step, which adds $W \cdot (n + \mathit{bonus}) - R$ to the opponent's damage, with base damage $n \ge 1$, Weakness $W$ of 1 or 2, and Resistance $R \le k$; a bonus of $2k$ or more makes that at least $k$ (\lean{dealtBy\_cap}). That damage is read only by tests against constants below $k$ and by the Knock Out test at an HP of at most $k$, so a bonus of $2k$ or more acts as $2k$ does. This needs tameness: no reduction subtracts from the bonus, and no \lean{drain} reads the damage done. \\
Damage reduction, both stores & $2k$ & Under \lean{Fit k} there is none, and tame cards add none. \\
\bottomrule
\end{tabularx}
\caption{What \lean{GameState.clamp k} caps, and why, for tame decks.}\label{tab:clamp}
\end{table}

\paragraph{Why the stores are capped at $2k$.} Resistance is a card constant that the
calculation subtracts after the bonus is added, so a stored bonus that is capped must stay large
enough to survive it. Capped at $k$ itself it does not: for every $k \ge 2$, bonuses of $k$ and
$k + 1$ are equal at the cap, but from base damage 1 against a Resistance of 3 they do $k - 2$
and $k - 1$, which differ below $k$ (\lean{Calc.cap\_at\_k\_fails}). This is a fact about the
damage function, given as the reason for the factor 2; it is not a game-level witness that a
clamp capping the stores at $k$ fails in play, and at $k = 2$ its Resistance of 3 exceeds what
\lean{Fit 2} allows. Hence \lean{HpLe} also bounds Resistance, and \lean{bound} includes the
largest Resistance. Weakness needs no larger cap: it doubles $n + b$ for a bonus $b$, and a cap
$K \ge k$ on $b$ does not change the result below $k$, since
$\min(2(n + b), k) = \min(2(n + \min(b, K)), k)$ (\lean{Calc.weak\_cap\_comm}).
The cap of $2k$ applies to every deck, whether or not its cards have Resistance.

\subsection{Simulation inside and between moves}\label{sec:sim}

Nothing is Knocked Out inside one \lean{Run}, and the tests see damage above HP, so the clamp
is not a simulation there.

\begin{theorem}[\lean{Reads.run\_clamp\_not\_sim}]\label{thm:notsim}
For every cap $k$, there are an effect and two states with the same clamp such that some
resolution from the first is matched by no resolution from the second up to the clamp.
\begin{lstlisting}
theorem run_clamp_not_sim (k : Nat) :
    ∃ e g₁ g₂ g₁', g₁.clamp k = g₂.clamp k ∧ Run e g₁ g₁' ∧
      ∀ g₂', Run e g₂ g₂' → g₁'.clamp k ≠ g₂'.clamp k
\end{lstlisting}
\end{theorem}
\noindent The witness has the opponent's 5-HP Active \pkmn{} carrying 5 damage counters in one
state and 6 in the other; the effect \lean{ite (atLeast opp 6) skip draw} draws a card in the
first and not in the second.

Inside one attack's effect, the player's own damage is tracked exactly and the opponent's up to
the cap $k$; Poison strength can still be capped at HP, since nothing in an effect reads it. A
damage step's amount may read the opponent's damage: if that damage is below $k$ it is the same
in both states, and if not, the opponent's damage after the step is at least $k$ in both,
whatever the amounts.

\begin{theorem}[\lean{Run.cap\_sim}, \lean{TRun.clamp\_sim}, \lean{Move.clamp\_sim},
\lean{Reach.clamp\_sim}]\label{thm:clamp}
For a tame effect whose test constants are below $k$, from a state with no damage reduction and
with every HP and Resistance at most $k$, states with the same \lean{rclamp k} have the same
effect resolutions, up to the same map. Under the same hypotheses, from a settled or finished
state, states
with the same clamp have the same Trainer-effect resolutions, up to the clamp. From a state
that is \lean{Fit k}, states with the same clamp have the same moves, and from start states
with the same clamp the same clamped states are reachable.
\begin{lstlisting}
theorem Run.cap_sim {k : Nat} (hk0 : 0 < k) {e : Effect} {g₁ g₁' : GameState}
    (h : Run e g₁ g₁') :
    e.tame = true → e.tmax < k → g₁.NoReduce → g₁.HpLe k → ∀ {g₂ : GameState},
      g₁.rclamp k = g₂.rclamp k →
      ∃ g₂', Run e g₂ g₂' ∧ g₁'.rclamp k = g₂'.rclamp k
theorem TRun.clamp_sim {k : Nat} (hk0 : 0 < k) {e : Effect} {g₁ g₁' : GameState}
    (h : TRun e g₁ g₁') :
    e.tame = true → e.tmax < k → g₁.NoReduce → g₁.HpLe k → ∀ {g₂ : GameState},
      g₁.clamp k = g₂.clamp k → (Settled g₁ ∨ g₁.winner ≠ none) →
        ∃ g₂', TRun e g₂ g₂' ∧ g₁'.clamp k = g₂'.clamp k
theorem Move.clamp_sim {k : Nat} {g₁ g₁' g₂ : GameState} (hF : g₁.Fit k) (h : Move g₁ g₁')
    (hc : g₁.clamp k = g₂.clamp k) : ∃ g₂', Move g₂ g₂' ∧ g₁'.clamp k = g₂'.clamp k
theorem Reach.clamp_sim {k : Nat} {g₀ h₀ g₁ : GameState} (hF : g₀.Fit k) (h : Reach g₀ g₁)
    (hc : g₀.clamp k = h₀.clamp k) : ∃ g₂, Reach h₀ g₂ ∧ g₁.clamp k = g₂.clamp k
\end{lstlisting}
\end{theorem}

\lean{Fit k} transfers from $g_1$ to any $g_2$ with the same clamp (by the lemma
\lean{GameState.Fit.of\_clamp}), so \lean{Move.clamp\_sim} gives both directions of a strong
bisimulation for \lean{Move} on \lean{Fit} states. The proof moves between the two maps at the
boundary of each move. Suppose $g_1$ and $g_2$ have the same clamp, $g_1$ is \lean{Fit k}, and
$g_1$ makes a move.
\begin{itemize}
\item If the move is one of the six Knock Out moves, its only read of a capped quantity is the
  test \lean{hp $\le$ damage} or its negation, which gives the same answer on the capped value,
  and its effect on the state (removing the \pkmn{}, moving a Prize card) commutes with the
  clamp.
\item Otherwise the move requires $g_1$ to be settled. Then every \pkmn{} in play is below its
  HP, the clamp does not change its damage, and so $g_1$ and $g_2$ have \emph{equal} damage
  everywhere; they differ at most in Poison strength above HP and in damage bonuses at or above
  $2k$, and $g_2$ is settled too. An attack's effect resolves by \lean{Run}, and
  \lean{Run.cap\_sim} matches it from $g_2$ with the player's own damage exact at every step,
  however far the effect drives it above HP and however often it reads it, and with the
  opponent's damage exact below $k$; Weakness and Resistance are read from the same cards in
  both states. At the end of the attack the Knock Out test sees the same
  answers, and the clamps agree. A Trainer effect resolves by \lean{TRun}; each of its steps
  starts settled (Theorem~\ref{thm:guarded}), so the same argument applies step by step, and
  the Knock Outs after each step are matched as in the first case. Checkup adds
  $1 + {}$\lean{poisonUp} to damage from a settled state, and the capped result depends on
  Poison strength only through its cap.
\end{itemize}
The argument depends on the Knock Out timing. Without \lean{Settled}, a move could start with
a \pkmn{} above its HP, \lean{heal} would read its exact damage, and two states with the same
clamp would part (with HP 3, damage 5 healed by 2 leaves 3, still at HP, whereas the clamped
damage 3 healed by 2 leaves 1). It depends on tameness (Section~\ref{sec:linger-fail}), and on
the load-bearing departures of Section~\ref{sec:departures}: a test on Poison strength against a
number above HP would observe the capped excess, and HP modifiers would move the cap.

\subsection{Finiteness}\label{sec:finite}

\begin{theorem}[\lean{Reach.finite}, \lean{Reach.clamp\_mem}, \lean{GameState.fit\_bound}]\label{thm:finite}
Let $g_0$ be a well-formed start state that is \lean{Fit} at $k = {}$\lean{bound g$_0$.allCards}.
Then the clamps at $k$ of all states reachable from $g_0$ lie in one finite list, and the clamp
is a bisimulation along reachable states: if $g_1$ is reachable and $g_2$ has the same clamp,
every move from either one is matched by a move from the other to a state with the same clamp.
The list half alone holds from every well-formed start. A start whose cards are all tame and
with no reduction in effect is \lean{Fit} at \lean{bound}.
\begin{lstlisting}
theorem Reach.finite {g₀ : GameState} (hw : WellFormed g₀) (hF : g₀.Fit (bound g₀.allCards)) :
    (∃ L : List GameState, ∀ g, Reach g₀ g → g.clamp (bound g₀.allCards) ∈ L) ∧
      ∀ g₁ g₂, Reach g₀ g₁ → g₁.clamp (bound g₀.allCards) = g₂.clamp (bound g₀.allCards) →
        (∀ g₁', Move g₁ g₁' → ∃ g₂', Move g₂ g₂' ∧
            g₁'.clamp (bound g₀.allCards) = g₂'.clamp (bound g₀.allCards)) ∧
        (∀ g₂', Move g₂ g₂' → ∃ g₁', Move g₁ g₁' ∧
            g₂'.clamp (bound g₀.allCards) = g₁'.clamp (bound g₀.allCards))
theorem Reach.clamp_mem {g₀ : GameState} (hw : WellFormed g₀) :
    ∃ L : List GameState, ∀ g, Reach g₀ g → g.clamp (bound g₀.allCards) ∈ L
theorem GameState.fit_bound {g₀ : GameState} (ht : ∀ x ∈ g₀.allCards, x.tame = true)
    (hn : g₀.NoReduce) : g₀.Fit (bound g₀.allCards)
\end{lstlisting}
\end{theorem}

So, for tame decks, the quotient of the reachable states by the clamp is a finite transition
system with the same moves, and any question about play whose answer depends only on the clamp
has the same answer in it: every question that does not read damage or Poison strength above
HP on an Active \pkmn{}, or a damage bonus above $2k$.

The statement's text does not show which clamp and which hypothesis it is about; the definitions
do. Spelled out, its hypotheses are a well-formed start whose cards are all tame and in which no
damage reduction is in effect; its cap is $k = {}$\lean{bound}, which exceeds every printed HP,
every printed Resistance and every test constant among the cards; and its clamp caps damage and
Poison strength on the two Active \pkmn{} at their printed HP and every stored bonus and
reduction at $2k$. Weakness and Resistance are allowed. Because the stores are capped at $2k$
rather than at $k$, the finite quotient is finer than one with the stores capped at $k$, and the
bisimulation claim identifies fewer states.

\lean{Reach.clamp\_mem} is the unconditional half. It is a statement about the clamp, not about
play: without a hypothesis that makes the clamp a bisimulation it says nothing about which
questions of play the finite list answers, and Theorem~\ref{thm:breaks} gives decks for which
play separates states with the same clamp. Its witness is explicit: the list
\lean{states g$_0$.allCards}, built from all lists of length at most $N$ over the $N$ starting
cards for each zone, the \pkmn{} cards among them as top cards, damage and Poison strength up to
\lean{bound} and every damage bonus and reduction up to twice \lean{bound}, the four rotation
conditions, and every
value of each flag and of each \pkmn{}'s markers. It is a crude over-approximation, it is never
computed, and its length is astronomically larger than the true reachable set. The proof uses
Theorem~\ref{thm:perm} (cards are only permuted) and \lean{Reach.wellFormed} (Benched \pkmn{}
are below HP and have no Special Condition). The cap on damage at move boundaries is the
printed HP; only the cap on bonuses must be large, at twice a number above the cards' test
constants and Resistance, because a bonus is read inside an attack, where tests see damage
exactly and Resistance is subtracted after it.

\subsection{Where the clamp fails}\label{sec:linger-fail}

\begin{theorem}[\lean{reduce\_breaks\_clamp}, \lean{drain\_breaks\_clamp}]\label{thm:breaks}
For every $k \ge 3$ there is a well-formed, settled start with no reduction in effect, every HP
and Resistance at most $k$ and every test constant below $k$, in which one card, an Item, adds a damage
reduction, and from which two states with the same \lean{clamp k} are reachable such that some
move from the first is matched by no move from the second up to \lean{clamp k}. The same holds
for every $k \ge 2$ with a start in which one card's attack uses \lean{drain} instead of a
reduction.
\begin{lstlisting}
theorem reduce_breaks_clamp (k : Nat) (hk : 3 ≤ k) :
    ∃ g₀, WellFormed g₀ ∧ Settled g₀ ∧ g₀.NoReduce ∧ g₀.HpLe k ∧
      (∀ x ∈ g₀.allCards, x.tmax < k) ∧
      ∃ g₁ g₂ g₁', Reach g₀ g₁ ∧ Reach g₀ g₂ ∧ g₁.clamp k = g₂.clamp k ∧ Move g₁ g₁' ∧
        ∀ g₂', Move g₂ g₂' → g₁'.clamp k ≠ g₂'.clamp k
theorem drain_breaks_clamp (k : Nat) (hk : 2 ≤ k) :
    ∃ g₀, WellFormed g₀ ∧ Settled g₀ ∧ g₀.NoReduce ∧ g₀.HpLe k ∧
      (∀ x ∈ g₀.allCards, x.tmax < k) ∧
      ∃ g₁ g₂ g₁', Reach g₀ g₁ ∧ Reach g₀ g₂ ∧ g₁.clamp k = g₂.clamp k ∧ Move g₁ g₁' ∧
        ∀ g₂', Move g₂ g₂' → g₁'.clamp k ≠ g₂'.clamp k
\end{lstlisting}
(Both in namespace \lean{Pokemon.Lingering}.)
\end{theorem}
\noindent The statements cover every cap that matters, at or above the deck's HP and test
constants; smaller caps are not considered, because \lean{Fit k} asks every printed HP to be at
most $k$ and the witnesses have 3-HP and 2-HP \pkmn{}. No card in either witness has Weakness or
Resistance. In the first witness the Item gives the opponent a reduction of $2k+1$ for this turn
and the player a bonus of $2k+1$ or $2k+2$, by choice. Both bonuses lie above the store cap
$2k$, so the two states have the same clamp. An attack with base damage 1 then does
$(1 + (2k+1)) - (2k+1) = 1$ or $(1 + (2k+2)) - (2k+1) = 2$ damage to a target with 3 HP, and
the clamps of the results differ. In the second witness the attacker, with 2 HP, has an attack
whose text puts $2k+3$ damage counters on itself, above every cap, and then makes its damage
step, \lean{drain (const 1)}, which heals it by the damage done, $1 + (2k+1)$ or $1 + (2k+2)$,
leaving 1 damage counter or none.

The damage done is the difference of two quantities that can each grow without limit, and
\lean{damage} exposes that difference to tests and to the Knock Out test; a cap on each
quantity separately cannot preserve the difference. The first witness uses only the
\lean{now} slot, so restricting the rest-of-game slot alone does not rescue the clamp. Each
theorem exhibits one deck outside the tame class, not every such deck (a deck whose reduction
card is never played is unaffected). Both witness attacks have base damage 1, so the rulebook's
stop at base damage 0 does not apply to them.

\subsection{Exact stores and faithful decks}\label{sec:faithful}

Both witnesses of Theorem~\ref{thm:breaks} reach two states whose stores differ only above the
cap. A clamp that keeps the stores exact does not identify them. \texttt{Exact.lean} defines it:
\begin{lstlisting}
def GameState.xclampAt (co cp : Cap) (g : GameState) : GameState :=
  { g with cur := { g.cur with active := g.cur.active.map (InPlay.clampWith co) },
           opp := { g.opp with active := g.opp.active.map (InPlay.clampWith cp) } }

def GameState.xclamp (g : GameState) : GameState := g.xclampAt .hp .hp
\end{lstlisting}
\noindent The \emph{exact-store clamp} \lean{xclamp} caps damage and Poison strength on the two
Active \pkmn{} at the printed HP and changes nothing else; it has no parameter. Every reader of
the stores sees equal values in two states with the same exact-store clamp: the damage step's
amount evaluates the same, Weakness and Resistance are read from the same cards and types, and
the stores are equal, so the step computes the same \lean{dealtBy}, however large, and
\lean{drain} heals the attacker by that amount; locks and prevention are flags in equal stores;
and \lean{endTurn} shifts equal stores. Inside an attack the relation \lean{xrclamp} keeps
damage exact on both Active \pkmn{}, since every attack starts settled, where the exact-store
clamp determines the damage. Between moves damage is read by Checkup only from settled states,
and by the Knock Out test at HP.

\begin{theorem}[\lean{GameState.xclamp\_bisim}]\label{thm:xbisim}
For every deck and any two states with the same exact-store clamp, reachable or not, every move
from either is matched by a move from the other to a state with the same exact-store clamp.
\begin{lstlisting}
theorem GameState.xclamp_bisim {g₁ g₂ : GameState} (hc : g₁.xclamp = g₂.xclamp) :
    (∀ g₁', Move g₁ g₁' → ∃ g₂', Move g₂ g₂' ∧ g₁'.xclamp = g₂'.xclamp) ∧
      (∀ g₂', Move g₂ g₂' → ∃ g₁', Move g₁ g₁' ∧ g₂'.xclamp = g₁'.xclamp)
\end{lstlisting}
\end{theorem}
\noindent No hypothesis on the deck is needed: damage reduction and \lean{drain} are allowed.
What the exact stores give up is a bound on the stores, and with it finiteness. On states whose
stores are bounded, finiteness returns (\texttt{Bounded.lean}):
\begin{lstlisting}
def GameState.BoundedStores (B : Nat) (g : GameState) : Prop :=
  StoreLe B g.cur.mods ∧ StoreLe B g.opp.mods

theorem Reach.xfinite_bounded {g₀ : GameState} (hw : WellFormed g₀) (B : Nat) :
    (∃ L : List GameState, ∀ g, Reach g₀ g → g.BoundedStores B → g.xclamp ∈ L) ∧
      ∀ g₁ g₂, Reach g₀ g₁ → g₁.xclamp = g₂.xclamp →
        (g₁.BoundedStores B → g₂.BoundedStores B) ∧
        (∀ g₁', Move g₁ g₁' → ∃ g₂', Move g₂ g₂' ∧ g₁'.xclamp = g₂'.xclamp) ∧
        (∀ g₂', Move g₂ g₂' → ∃ g₁', Move g₁ g₁' ∧ g₂'.xclamp = g₁'.xclamp)
\end{lstlisting}
\noindent \lean{BoundedStores B} says that every slot of both stores has a bonus and a
reduction of at most $B$. The bounded part is not closed under moves, so on its own this gives
finitely many classes and a bisimulation that may leave them, not a finite quotient. A condition
on the deck keeps every reachable state in the bounded part (\texttt{Faithful.lean}):
\begin{lstlisting}
def Effect.loopFree : Effect → Bool
  | .ite _ a b => a.loopFree && b.loopFree
  | .seq a b => a.loopFree && b.loopFree
  | .choice a b => a.loopFree && b.loopFree
  | .iter a => !a.writes
  | _ => true

def GameState.LoopFree (g : GameState) : Prop := ∀ x ∈ g.allCards, x.loopFree = true

def GameState.GXOnly (g : GameState) : Prop :=
  ∀ p : PokemonCard, Card.pokemon p ∈ g.allCards → ∀ atk ∈ p.attacks, atk.gx = false →
    atk.effect.amtAt .always = 0

def GameState.Faithful (g : GameState) : Prop := g.LoopFree ∧ g.GXOnly
\end{lstlisting}
\noindent Here \lean{a.writes} says that \lean{a} contains a damage bonus or reduction, and
\lean{amtAt .always} is the largest rest-of-game amount one resolution of an effect writes. A
start is \emph{faithful} if no card puts a damage bonus or reduction under \lean{iter}
(\lean{LoopFree}) and every attack that writes a rest-of-game bonus or reduction is a GX attack
(\lean{GXOnly}). Locks and prevention may appear anywhere, and Item and Supporter cards may
write rest-of-game amounts.

\begin{theorem}[\lean{Faithful.reach\_boundedStores}, \lean{Faithful.xfinite}]\label{thm:xfinite}
From a faithful start, every reachable state has stores bounded by a number computed from the
start state. From a well-formed faithful start, the exact-store clamps of all reachable states
lie in one finite list, and the exact-store clamp is a two-way bisimulation along reachable
states.
\begin{lstlisting}
def storeBound (g₀ : GameState) : Nat :=
  g₀.storeTotal + cardsAmt g₀.allCards + 3 * cardsAtkAmt g₀.allCards

theorem reach_boundedStores {g₀ g : GameState} (hF : g₀.Faithful) (h : Reach g₀ g) :
    g.BoundedStores (storeBound g₀)
theorem xfinite {g₀ : GameState} (hw : WellFormed g₀) (hF : g₀.Faithful) :
    (∃ L : List GameState, ∀ g, Reach g₀ g → g.xclamp ∈ L) ∧
      ∀ g₁ g₂, Reach g₀ g₁ → g₁.xclamp = g₂.xclamp →
        (∀ g₁', Move g₁ g₁' → ∃ g₂', Move g₂ g₂' ∧ g₁'.xclamp = g₂'.xclamp) ∧
        (∀ g₂', Move g₂ g₂' → ∃ g₁', Move g₁ g₁' ∧ g₂'.xclamp = g₁'.xclamp)
\end{lstlisting}
(All three are in namespace \lean{Pokemon.Faithful}.)
\end{theorem}
\noindent The bound \lean{storeBound} adds the start stores, the amounts of every Trainer card
in the game, and three times the largest amount of one attack. It rests on three facts of the
model. No effect or move returns a card from a discard pile, so each Item and Supporter is
played at most once per game. At most one attack is made per turn. And carried amounts expire
within three turn boundaries, while rest-of-game amounts written by attacks come only from GX
attacks, at most one per player per game. Faithful decks may use damage reduction (Iron
Defender), \lean{drain}, carried bonuses and reductions in attacks and Trainer cards, and
rest-of-game modifiers in GX attacks and in Trainer cards.

\paragraph{Limits of the faithful theorem.} The store bound rests on three absences in the
model: no effect returns a card from a discard pile, at most one attack is made per turn, and
there are no extra turns. Each exists on printed cards (Item- and Supporter-recycling cards,
``attack twice'', ``take another turn''), so adding any of them to the model reopens the bound.
The faithful class also does not cover copy attacks (for instance copying a GX attack);
static modifiers and modifiers from Tools, Abilities or Special Energy, since every modifier in
the model is written by an effect when it resolves; reductions restricted by type (the printed
Iron Defender protects Metal \pkmn{} only, whereas the model's reduction protects all of a
player's \pkmn{}); or effects that lock GX attacks. Rest-of-game modifiers from VSTAR powers,
Abilities or Stadiums are absent from the model. Today's Standard pool has card recycling and
attacking twice, though no extra turns, so the theorem does not yet reach it
(Section~\ref{sec:standard}).

\paragraph{Which theorem applies.} Theorem~\ref{thm:finite} needs tame cards and no reduction
in effect at the start; it allows bonus writes under \lean{iter} and rest-of-game bonuses in
attacks that are not GX attacks, and its quotient is coarser, since the clamp caps stored
bonuses at $2k$. Theorem~\ref{thm:xfinite} allows damage reduction, \lean{drain}, carried modifiers and
modifiers on Item and Supporter cards; it forbids amounts written under \lean{iter} and
rest-of-game amounts written by attacks that are not GX attacks, and its quotient is finer,
since the stores are kept exact. Neither class contains the other, and for decks in both, both
theorems apply. Section~\ref{sec:counter-sharp} shows that the GX clause cannot be dropped.

\subsection{Consequences and scope}\label{sec:finite-scope}

\begin{corollary}[argued, not mechanized]\label{cor:counter}
Fix a well-formed start state $g_0$ of the model that is either \lean{Fit} at
\lean{bound g$_0$.allCards} or faithful. Say that $g_0$ \emph{simulates} a machine $M$ if there
is an injective map from the configurations reachable by $M$ to classes of states reachable
from $g_0$, under the clamp in the first case and the exact-store clamp in the second. Then $M$
has finitely many reachable configurations. In particular, no fixed tame or faithful deck whose
card text is expressible in \lean{Effect} simulates a two-counter machine with infinitely many
reachable configurations.
\end{corollary}
\begin{proof}
By Theorem~\ref{thm:finite} or Theorem~\ref{thm:xfinite}, the classes reachable from $g_0$ lie
in a finite list.
\end{proof}
\noindent The corollary is immediate; its content lies in Theorems~\ref{thm:finite}
and~\ref{thm:xfinite}, which say that the game reachable from $g_0$ is bisimilar to a finite
transition system. The notion of simulation is deliberately weak, so the corollary excludes any
encoding that needs infinitely many distinguishable configurations from one tame or faithful
start. The construction of Section~\ref{sec:linger} uses, for every machine whose first
instruction is not \lean{halt}, a deck that is neither, and so lies outside it.

\begin{remark}[decidability]\label{rem:decide}
We conjecture that reachability is decidable for tame decks, and likewise for faithful decks
with the exact-store quotient; this is not formalized. The
finite list bounds the clamped reachable set, so a decision procedure needs only a computable
clamped successor relation and decidable equality on states. For \lean{TRun} a saturation
argument should give the successors, because every step starts settled. For \lean{Run} it is
less simple, since inside one attack's effect \lean{iter} can manipulate an unbounded damage
value that tests read exactly. Over the start states built in Section~\ref{sec:linger}, whose
decks are not tame, no uniform procedure exists (Section~\ref{sec:counter-what}).
\end{remark}

\paragraph{Scope.} The theorems of this section concern the states between moves, reached from
well-formed starts, for decks whose card text is expressible in \lean{Effect}. Inside one
attack's effect, damage is unbounded and readable (Theorems~\ref{thm:unbounded}
and~\ref{thm:notsim}). The finite list depends on $g_0$ and gives no bound uniform in the deck.
Coin flips and shuffles are nondeterministic choices, so the theorems are silent about
probability. \lean{Reach.perm} conserves the cards of both players together, not each player's
cards on that player's side. Two facts are argued and not proved: that each primitive step
damages at most one Active \pkmn{}, so that simultaneous Knock Outs cannot arise inside a
Trainer effect (the documentation of \texttt{Rules.lean} relies on this), and that
\lean{Resolve} always has a derivation from a well-formed state. The model's Trainer timing
follows the Advanced Player's Rulebook, not an English rule, and only up to the caveats of
Section~\ref{sec:departures}. Printed decks are not all tame, and the model
omits HP modifiers and tests on Poison strength, on which the proof depends; it keeps markers
as flags under an assumption about the card pool, and the faithful bound rests on the model's
lack of card recycling, repeated attacks and extra turns. Within this scope, a tame deck from a
well-formed start with no reduction in effect, and a faithful deck from a well-formed start,
each have a finite clamped game with the same moves as the game itself.

%======================================================================
\section{The universal side}\label{sec:linger}
%======================================================================

This section shows that the same model, for decks that no printed card supplies, simulates
every two-counter machine. The construction is in \texttt{Pokemon/Counter.lean}. It uses the
existing effect language and changes no rule; what makes it possible is card text more general
than any printed card, listed in Section~\ref{sec:counter-generality}. In this section all Lean
text is in namespace \lean{Pokemon.Counter}.

\subsection{Two-counter machines}\label{sec:counter-machines}

A two-counter (Minsky) machine \cite{minsky} is a finite program over two registers. The
instruction \lean{inc r q} increments register \lean{r} and goes to state \lean{q};
\lean{jz r $q_1$ $q_2$} goes to $q_1$ if register \lean{r} is zero and otherwise decrements it
and goes to $q_2$; \lean{halt} halts. Every state at or beyond the end of the program halts.
\begin{lstlisting}
inductive Reg where
  | one | two

inductive Instr where
  | inc (r : Reg) (q : Nat)
  | jz (r : Reg) (q₁ q₂ : Nat)
  | halt

structure Machine where
  prog : List Instr

structure Config where
  pc : Nat
  a : Nat
  b : Nat

def instr (M : Machine) (q : Nat) : Instr := M.prog.getD q .halt

inductive Step (M : Machine) : Config → Config → Prop where
  | incOne {pc a b q : Nat} (h : M.instr pc = .inc .one q) : Step M ⟨pc, a, b⟩ ⟨q, a + 1, b⟩
  | incTwo {pc a b q : Nat} (h : M.instr pc = .inc .two q) : Step M ⟨pc, a, b⟩ ⟨q, a, b + 1⟩
  | zeroOne {pc b q₁ q₂ : Nat} (h : M.instr pc = .jz .one q₁ q₂) :
      Step M ⟨pc, 0, b⟩ ⟨q₁, 0, b⟩
  | decOne {pc a b q₁ q₂ : Nat} (h : M.instr pc = .jz .one q₁ q₂) :
      Step M ⟨pc, a + 1, b⟩ ⟨q₂, a, b⟩
  | zeroTwo {pc a q₁ q₂ : Nat} (h : M.instr pc = .jz .two q₁ q₂) :
      Step M ⟨pc, a, 0⟩ ⟨q₁, a, 0⟩
  | decTwo {pc a b q₁ q₂ : Nat} (h : M.instr pc = .jz .two q₁ q₂) :
      Step M ⟨pc, a, b + 1⟩ ⟨q₂, a, b⟩

inductive Reaches (M : Machine) (c₀ : Config) : Config → Prop where
  | refl : Reaches M c₀ c₀
  | step {c c' : Config} : Reaches M c₀ c → Step M c c' → Reaches M c₀ c'

def Halted (M : Machine) (c : Config) : Prop := M.instr c.pc = .halt

def Halts (M : Machine) (c₀ : Config) : Prop := ∃ c, M.Reaches c₀ c ∧ M.Halted c

def init (a b : Nat) : Config := ⟨0, a, b⟩
\end{lstlisting}
\noindent (\lean{instr} through \lean{init} are in namespace \lean{Machine}.) The step relation
is deterministic, by the lemma \lean{Machine.Step.det}, and no step leaves a halted
configuration, by \lean{Machine.Step.not\_halted}.

\subsection{The construction}\label{sec:linger-counter}

The start state \lean{start M a b} uses two \pkmn{} cards built for $M$, the \emph{Clock}
(\lean{clock M}) and the \emph{Tape} (\lean{tape M}), each with one attack of cost nothing and
base damage 0. Let \lean{span} be one more than the largest state any instruction can go to (at
least 1), and \lean{W M = 2 * span}.
\begin{lstlisting}
def clock (M : Machine) : PokemonCard :=
  ⟨"Clock", .basic, none, 3 * V M + 1, .colorless, [atkA M], 0, none, none⟩

def tape (M : Machine) : PokemonCard := ⟨"Tape", .basic, none, 2, .colorless, [atkB M], 0, none, none⟩

def start (M : Machine) (a b : Nat) : GameState where
  cur := ⟨[], [.energy .grass], [], [.energy .grass], [], some (onPlay (clock M)), [],
    false, false, false, false, false, storeOf 0 b, true⟩
  opp := ⟨[.energy .grass], [], [], [.energy .grass], [], some (onPlay (tape M)), [],
    false, false, false, false, false, storeOf 0 a, true⟩
  curIsFirst := true
  firstTurn := false
  attacked := false
  between := false
  winner := none
  stadium := none
\end{lstlisting}
\noindent Here \lean{V M} is \lean{span}, \lean{onPlay} puts a Basic \pkmn{} into play with no
damage, Energy or markers, and \lean{storeOf b r} is a store whose only modifiers are a
rest-of-game damage bonus \lean{b} and reduction \lean{r}. The first player, to move in the
turn proper, has the Clock Active, one Energy card in hand, an empty deck and one Prize card;
the second player has the Tape Active, one Energy card in the deck, an empty hand and one Prize
card. Neither has a Bench, both have had a turn, and neither \pkmn{} has Weakness or Resistance.
Since both attacks have base damage 0, each attack's script is its effect. The construction has
six parts.
\begin{itemize}
\item \emph{Registers.} The rest-of-game slot never expires. Register one is the second player's
  rest-of-game damage reduction minus the first player's rest-of-game damage bonus; register two
  is the first player's reduction minus the second player's bonus. The Clock's effect
  \lean{linger .opp .always (.reduce 1)} increments register one and
  \lean{linger .own .always (.reduce 1)} increments register two; \lean{linger .own .always
  (.bonus 1)} decrements register one when the Clock uses it and register two when the Tape
  does. The inputs $a$ and $b$ are the start reductions.
\item \emph{Zero test.} The damage step \lean{damage (const 1)} does
  $1 + \mathit{bonus} - \mathit{reduction}$, truncated at 0. Read in the right direction, that is
  1 if the register is zero and 0 otherwise, and \lean{ite (exactly opp v)} then reads the result
  exactly. The Clock's attack tests register one, the Tape's register two. Each attack has its
  damage step in one branch only, so it makes the damage calculation at most once
  (\lean{atkA\_wf}, \lean{atkB\_wf}). This is the only place where a test reads damage placed
  by an earlier step in the same attack.
\item \emph{Program counter.} The machine state is the damage on the Clock. The Clock sets it
  by \lean{healAll} followed by \lean{selfDamage}, and reads it by \lean{exactly own q}; the
  Tape reads it by \lean{exactly opp q}. After running an instruction the Clock writes
  $W + q'$ for the next state $q'$. The Tape can only add damage to the Clock (by
  \lean{counters}), so after its test it adds damage counters until the Clock shows $W + q'$.
  The Clock's attack first normalises $W + q$ to $q$, then runs the instruction of state $q$
  unless it is a test of register two, which it leaves to the Tape. The Tape acts only when the
  Clock shows the normalised form $q$, so each round runs exactly one machine step. The Clock's
  HP, $3 \cdot \mathit{span} + 1$, exceeds every value the program counter takes, so printed HP
  grows with the machine.
\item \emph{Clean-up.} The Clock's zero test damages the Tape, which has 2 HP, by at most 1.
  The Tape's attack heals it (\lean{healAll}) first, and the Clock's attack runs nothing while
  the Tape is damaged. No \pkmn{} ever reaches its HP, so there is no Knock Out and no Prize
  card is taken.
\item \emph{Decks.} Each attack begins with \lean{handToBottom}, which returns the card just
  drawn to the deck, so on the intended line the decks never run out. A player who declines to
  attack, or who attaches the Energy card instead of keeping it in hand, decks out at their
  next draw, and the game ends (\lean{deckOut}) with a winner and without a halt marker.
\item \emph{Halt marker.} In a halting state the Clock's attack puts an Imprison marker on the
  Clock itself (\lean{mark .own .imprison}).
\end{itemize}
One machine step takes six moves: the first player attacks with the Clock and ends the turn;
the second player draws, attacks with the Tape and ends the turn; the first player draws. The
encoding of a configuration and the halt marker are:
\begin{lstlisting}
def Encodes (M : Machine) (c : Config) (g : GameState) : Prop :=
  ∃ s : Abs, s.cfg = c ∧ Good M s ∧ s.rB = 0 ∧ s.halted = false ∧
    Shape (pkA M s false true) (pkB M s true false) g ∧ g.between = false ∧
    g.attacked = false ∧ g.firstTurn = false ∧ g.winner = none

def HaltMarked (g : GameState) : Prop :=
  ∃ p, (g.cur.active = some p ∨ g.opp.active = some p) ∧ p.marks.imprison = true
\end{lstlisting}
\noindent Here \lean{Abs} is an abstract state: the machine state and registers, the two
players' rest-of-game bonuses, the damage on each \pkmn{}, and whether the Clock is marked.
\lean{Good M s} says that the program counter is $q$ or $W + q$ with $q$ below \lean{span},
that the Tape has at most one damage counter, and that the marker appears only in a halting
state. \lean{Shape} says that the game state has the shape the abstract state describes: one
Active \pkmn{} per player with the stated damage and markers and no Special Condition, no
Bench, only Energy cards in deck and hand, and only rest-of-game modifiers. \lean{Encodes}
holds at the start of the first player's turn proper, before attacking.

\subsection{The theorem}\label{sec:counter-theorem}

\begin{theorem}[\lean{Counter.universal}]\label{thm:universal}
For every two-counter machine $M$ and inputs $a$, $b$, the start state \lean{start M a b} is
well formed; it encodes the initial configuration; every configuration that $M$ reaches is
encoded by a reachable state; and a state with the halt marker is reachable if and only if $M$
halts.
\begin{lstlisting}
theorem universal (M : Machine) (a b : Nat) :
    WellFormed (start M a b) ∧ Encodes M (Machine.init a b) (start M a b) ∧
      (∀ c, M.Reaches (Machine.init a b) c → ∃ g, Reach (start M a b) g ∧ Encodes M c g) ∧
      ((∃ g, Reach (start M a b) g ∧ HaltMarked g) ↔ M.Halts (Machine.init a b))
\end{lstlisting}
It is assembled from these, among others:
\begin{lstlisting}
theorem simulate {M : Machine} {c c' : Config} {g : GameState} (hst : M.Step c c')
    (he : Encodes M c g) : ∃ g', Reach g g' ∧ Encodes M c' g'
theorem Encodes.inj {M : Machine} {c c' : Config} {g : GameState} (h : Encodes M c g)
    (h' : Encodes M c' g) : c = c'
theorem haltMarked_of_halts {M : Machine} {a b : Nat} (h : M.Halts (Machine.init a b)) :
    ∃ g, Reach (start M a b) g ∧ HaltMarked g
theorem halts_of_haltMarked {M : Machine} {a b : Nat} {g : GameState}
    (h : Reach (start M a b) g) (hm : HaltMarked g) : M.Halts (Machine.init a b)
\end{lstlisting}
\end{theorem}

\lean{simulate} is the forward simulation, one machine step by six moves, and
\lean{Encodes.inj} says that a state encodes at most one configuration. The two directions of
the equivalence differ in kind. Completeness, \lean{haltMarked\_of\_halts}, is
\emph{cooperative reachability}: if $M$ halts, some sequence of legal moves of the two players
reaches a state with the halt marker. It says nothing about whether either player would choose
those moves. Soundness, \lean{halts\_of\_haltMarked}, holds over all play: if any sequence of
legal moves reaches the halt marker, then $M$ halts.

\paragraph{Soundness over all play.} Soundness must cover every move either player may make,
not only the intended line. It rests on an invariant of all play from the start:
\begin{lstlisting}
def Inv (M : Machine) (c₀ : Config) (g : GameState) : Prop :=
  ∃ s, Good M s ∧ M.Reaches c₀ s.cfg ∧
    (Shape (pkA M s false false) (pkB M s false false) g ∨
      Shape (pkB M s false false) (pkA M s false false) g)

theorem move_inv {M : Machine} {c₀ : Config} {g g' : GameState} (hm : Move g g')
    (hi : Inv M c₀ g) : Inv M c₀ g'
\end{lstlisting}
\noindent \lean{Inv} says that the state has the shape of a \lean{Good} abstract state whose
configuration $M$ reaches, with either player to move, and \lean{move\_inv} says that every
legal move preserves it. Its proof is a case split over all seventeen constructors of
\lean{Move}, which Lean forces to be complete, and over every resolution of the attacks'
effects. Attacks run the machine one step or do nothing. Attaching Energy, drawing and ending
the turn change nothing abstract. Decking out ends the game. No other move is legal: hands hold
only Energy, Benches are empty, and no \pkmn{} is Confused or at its HP. A player who deviates
from the intended line decks out at their next draw, and the invariant covers \lean{deckOut}.
Finally, in a state satisfying \lean{Inv} the halt marker appears only when the abstract state
is halted (\lean{Inv.halts}), and the abstract configuration is one that $M$ reaches, so $M$
halts.

\subsection{The generality the construction uses}\label{sec:counter-generality}

The decks of the construction are expressible in the model but supplied by no printed card.
The module docstring of \texttt{Counter.lean} lists the departures from printed cards and from
legal play on which the construction relies:
\begin{enumerate}
\item rest-of-game damage bonuses and reductions that add up without limit, placed by an
  attack's effect that is used every turn (\lean{linger} is a general primitive), including an
  attack that grants the \emph{opponent} a rest-of-game damage reduction
  (\lean{linger .opp .always (reduce 1)}), whereas printed cards give rest-of-game modifiers
  only through once-per-game attacks (Section~\ref{sec:linger-printed}); neither attack is a GX
  attack, so the once-per-game rule does not restrict them, and this is exactly the clause
  \lean{GXOnly} of faithfulness that the construction violates (Section~\ref{sec:counter-sharp});
\item the damage placed by an attack read by a test later in the same attack (\lean{ite (exactly
  opp v)} after the damage step \lean{damage (const 1)}); the rulebook places damage and ``then
  do[es] all other effects'' \cite[p.~20]{rb-151}, so reading it afterwards is in order, and
  what is general is that the effect tests the exact damage and branches on it, and that the
  damage step sits inside a branch chosen by earlier tests of the same attack;
\item arbitrary printed HP: the Clock's HP, $3 \cdot \mathit{span} + 1$, grows with the machine;
\item a halt marker that the Clock places on itself (\lean{mark .own .imprison}), whereas printed
  markers target the opponent, which affects only the observable, not the simulation;
\item a hand-built start, in which each player has an Active \pkmn{}, one Energy card and one
  Prize card, with no set-up; it is not shown to be reachable from a legal set-up, and padding
  the decks to legal size is not proved harmless;
\item the model's other generalisations of card text: \lean{ite}, \lean{healAll},
  \lean{selfDamage}, \lean{counters}, \lean{mark} and \lean{handToBottom} as general
  primitives combined in one attack.
\end{enumerate}
The construction runs under the model's turn rules: each player attacks once and the turn ends.
It does not rely on a bonus being added to an attack with base damage 0: both attacks have
printed base damage 0, and the zero test is the damage step \lean{damage (const 1)}, with base
damage 1. Both attacks cost nothing, so typed costs do not restrict them. Each player has one
\pkmn{} and no Bench, so it does not use modifiers that belong to a player rather than to one
\pkmn{}. The Clock and the Tape have neither Weakness nor Resistance, each attack makes the
damage calculation at most once, and the construction does not use \lean{drain}.

Unlimited rest-of-game stacking is therefore one ingredient of one construction, used together
with the other five. The construction is a theorem about decks that no printed card supplies,
and it does not show that the printed game is Turing complete;
Section~\ref{sec:linger-printed} explains why the printed cards do not run it.

\subsection{The GX clause cannot be dropped}\label{sec:counter-sharp}

\texttt{Dichotomy.lean} relates the construction to the faithful theorem
(Theorem~\ref{thm:xfinite}). All three statements are in namespace \lean{Pokemon.Faithful}.

\begin{theorem}[\lean{counter\_loopFree}, \lean{counter\_not\_gxOnly} and
\lean{xclamp\_not\_finite}]\label{thm:sharp}
Every two-counter start satisfies the first clause of faithfulness. It violates the second
whenever the machine's first instruction is not \lean{halt}. And for \lean{start countUp 0 0}
the exact-store clamps of the reachable states lie in no finite list.
\begin{lstlisting}
theorem counter_loopFree (M : Machine) (a b : Nat) : (start M a b).LoopFree
theorem counter_not_gxOnly (M : Machine) (a b : Nat) (h : M.instr 0 ≠ .halt) :
    ¬ (start M a b).GXOnly
theorem xclamp_not_finite :
    ¬ ∃ L : List GameState, ∀ g, Reach (start countUp 0 0) g → g.xclamp ∈ L
\end{lstlisting}
\end{theorem}
\noindent The Clock's or the Tape's attack, neither a GX attack, writes a rest-of-game bonus or
reduction: an increment writes a reduction and a decrement a bonus. The theorem
\lean{Faithful.dichotomy} packages Theorem~\ref{thm:xfinite} with these facts and the halting
equivalence of Theorem~\ref{thm:universal}. Despite its Lean name, it is a sharpness result for
the clause \lean{GXOnly}, not a partition of decks: faithful decks are finite up to a two-way
bisimulation, and some decks that violate only \lean{GXOnly} simulate every two-counter
machine, so that clause cannot be dropped from Theorem~\ref{thm:xfinite}. Nothing is proved
about other decks that are not faithful, and some of them, tame decks that violate
\lean{GXOnly}, are finite by Theorem~\ref{thm:finite}. A machine whose first instruction is
\lean{halt} halts at once, and nothing is claimed about its start. Like
Theorem~\ref{thm:notfiniteupto} below, \lean{xclamp\_not\_finite} concerns representation:
\lean{countUp} never reads its register.

\subsection{What follows}\label{sec:counter-what}

\paragraph{Undecidability.} The map $M \mapsto{}$\lean{start M a b} is an explicit definition,
and by Theorem~\ref{thm:universal} it reduces the halting problem for two-counter machines to
reachability of the halt marker. With Minsky's theorem that this halting problem is undecidable
\cite{minsky}, reachability of the halt marker, as a problem over the family of start states
the reduction builds, is undecidable. The reduction is machine-checked; the appeal to Minsky's
theorem is not, and undecidability is not stated as a Lean theorem. It follows that no uniform,
computable finite abstraction answers the halt question over this family, since one would
decide the halting problem. For any single machine, a finite quotient that respects the halt
question exists, non-constructively, as \texttt{Counter.lean} notes.

\paragraph{Representation.} The file also proves that, for one witness deck, the reachable
states are not finite up to any equivalence fine enough to recover the encoded configuration.
The witness machine \lean{countUp} has one instruction, \lean{inc .one 0}: it counts up for
ever.
\begin{theorem}[\lean{Counter.not\_finite\_up\_to}]\label{thm:notfiniteupto}
Let $E$ be an equivalence on game states under which encodings of different configurations of
\lean{countUp} are never equivalent. Then no finite list contains, up to $E$, every state
reachable from \lean{start countUp 0 0}.
\begin{lstlisting}
def countUp : Machine := ⟨[.inc .one 0]⟩

theorem not_finite_up_to (E : GameState → GameState → Prop) (hE : Equivalence E)
    (hsep : ∀ c c' g g', Encodes countUp c g → Encodes countUp c' g' → E g g' → c = c') :
    ¬ ∃ L : List GameState, ∀ g, Reach (start countUp 0 0) g → ∃ h ∈ L, E g h
\end{lstlisting}
\end{theorem}
\noindent The case where $E$ is equality is \lean{Counter.not\_finite}. The hypothesis
\lean{hsep} asks $E$ to recover the encoded configuration whether or not play can observe it,
and \lean{countUp} never reads its register, so the extra states differ only in a damage
reduction that nothing observes. The theorem shows that representations of configurations are
infinite, and \lean{xclamp\_not\_finite} shows the same for the exact-store clamp; whether a
finite bisimulation quotient exists for this deck is open. The behavioural
counterpart of the finiteness theorems is Theorem~\ref{thm:universal} with Minsky's theorem.

%======================================================================
\section{The printed game}\label{sec:printed}
%======================================================================

The theorems concern the model. This section asks how far they reach the printed game: what
the card text contains, how much of the rules the model covers and what today's Standard pool
contains, where in the printed game an unbounded store could hide, what remains open, and what
we conjecture.

\subsection{The survey}\label{sec:survey}

\paragraph{Data.} We use the \texttt{pokemon-tcg-data} dataset behind pokemontcg.io
\cite{ptcgdata}, at commit \texttt{90f27d0} (dated 2026-11-08), fetched on 2026-11-21 by a
single shallow clone. It holds 20,776 English card records in 177 set files; card text is in
the ability, attack and rules fields. Japanese-only cards are absent.

\paragraph{Method.} By \emph{the survey} we mean six searches, each recorded in a survey file.
The \emph{literature search} is reported in Section~\ref{sec:related} and the \emph{rules
audit} in Section~\ref{sec:game} and Appendix~\ref{app:inventory}. The \emph{store search} and
the \emph{control-flow search} used the LimitlessTCG full-text search \cite{limitless}, which
shows current English text with errata applied. Three stages ran over the dataset. The
\emph{regex scan} applied regular expressions for unbounded HP, Knock Out prevention, game-long
counts, unbounded markers and cards from outside the game to the ability, attack and rules
text. The \emph{census} removed reprints with identical text, split the remaining text into
36,239 sentences, reduced each sentence to a template by replacing numbers, the card's own
name, its own attack and ability names, and type words by placeholders, and assigned the
sentences to 26 effect classes by keyword rules; we then read by eye every distinct sentence
that the risk-class patterns selected. The \emph{targeted reading} read every sentence about
lingering effects, markers, HP and the Bench. Appendix~\ref{app:survey} gives the statistics.
A regular expression finds only the wordings it describes, so ``none found'' below means none
found by the stated patterns.

\paragraph{Findings: plainly bounded.} Every HP modifier is a fixed amount from a card in play,
a bonus scaled by a count of cards or Prizes (for example ``+$N$ HP for each [type] Energy
attached'', or Brambleghast's ``+50 HP for each Prize card your opponent has taken''), or a
conditional fixed maximum; no attack or once-per-turn effect adds HP persistently and
repeatably, and the largest printed HP in the dataset is 380. No card says that a \pkmn{}
cannot be Knocked Out while damage keeps accumulating: every prevention sets the remaining HP
to 10, moves the \pkmn{} out of play, or changes only where the cards go. ``This game'' text
consists of once-per-game flags, lasting modifiers from once-per-game GX attacks, win or loss
clauses, and face-up Prize flags; no card reads a count of events over the whole game. No card
brings a card in from outside the game. Only Croagunk's Poison Up says ``This effect can be
applied more than once'', and the resulting Poison damage per Checkup is bounded because the
Poisoned \pkmn{} is Knocked Out once its damage reaches HP (Section~\ref{sec:hide}). Every count
read by card text (``for each'', ``number of'', ``exactly'') is bounded by cards, by HP, by the
Prize count, or by a printed constant.\label{sec:findings}

\paragraph{Findings: not plainly bounded.} Four things need a modelling decision rather than a
bound. First, no card caps how many Imprison or Shock-wave markers (Gardevoir ex $\delta$,
Tyranitar ex $\delta$, EX Dragon Frontiers) may sit on one \pkmn{}, and the store search found
no ruling either way, so the count may be unbounded; but every card that reads them asks only
whether a \pkmn{} has \emph{any}, and every removal removes all of them, so each is an unbounded
tally that card text reads only by a yes-or-no test. Second, 83 sentences say ``as often as you
like'' or ``as many times as you like'', and 194 match the census's flip-until-tails pattern;
each gives an unbounded number of repetitions, each of which moves cards or damage counters
within a finite set of states, so they add time, not memory. Third, two Ditto cards take their
HP from the opposing Active \pkmn{}, which is circular when both Active \pkmn{} are such Dittos.
Fourth, Rock-Paper-Scissors, guessing games about printed card data and physical actions (``sing
a song'', ``throw the card'') are not functions of the game state, and a model must treat them
as nondeterminism or exclude them.

\paragraph{The claim.} The survey found no unbounded store that card text reads beyond a
yes-or-no test. It leaves one case open: a damage modifier for the current turn is bounded by
the copies of its card times a printed constant only if no card-recycling loop replays the card
within the turn (Section~\ref{sec:linger-printed}).

\paragraph{Limits.}\label{sec:limits} The survey covers English text only; a Japanese-only card
that reads a marker count would change the marker finding. Each by-eye list is as complete as
the regular expression that selected it, and the weakest is the quantity-reading filter, which
may miss a quantity read implicitly. About 10\% of a by-eye sample of the census's class labels
were disputed, so the safety-relevant classes were read in full rather than trusted. The current
rulebook was not read: the rulebook on hand is the 2023 edition, and the later sources were read
from a third-party copy or through a summary (Section~\ref{sec:game}). The current TCG
Tournament Handbook was not read either, and some marker rulings are secondhand. Rules-level
questions, such as whether markers survive evolution, are not settled by card text.

\subsection{Rules coverage and today's Standard pool}\label{sec:standard}

Two further searches compare the model with the game as it stands today. Their findings are
surveyed in the sense of Section~\ref{sec:status}; counts are of distinct cards (name and full
text), found by regular expression, with the precision checks the search files record.

\paragraph{Rules coverage.} A catalogue of the game's rules and mechanisms has 124 entries, each
with its status in the model. At the time of the catalogue, 25 were faithful, 31 simplified and
68 absent. The catalogue predates two stages of the development: M1, the turn structure, and
M4, attack costs, Weakness and Resistance and the damage calculation. M1 settles the entries for
the end of the turn at an attack, the starting player's Supporter ban, the stop at base damage
0, a new Poison replacing the old one, and the nine Basic Energy types; it corrects the
per-player evolution ban, an entry that stays simplified because evolution reaches only the
Active \pkmn{}, and it makes the Stadium one shared slot, an entry that stays simplified because
no move plays a Stadium. M4 settles typed attack costs and the order of the damage calculation,
and adds Weakness and Resistance in their current forms. The catalogue also records the
once-per-game GX rule as unused, which the model has. We have not recounted the totals.

\paragraph{Today's Standard pool.} No rulebook defines the formats. We take Standard to be the
cards with regulation mark H, I or J, together with every Basic Energy; this rests on
search-result text, not on an official source, and is an assumption. The dataset's own Standard
flag is out of date. The pool so defined has 3,402 card records, 2,164 distinct cards and 1,470
names.

Several mechanisms that bear on the theorems are absent from it: GX attacks and VSTAR Powers,
rest-of-game modifiers (the only two ``rest of the game'' texts keep a Prize card face up:
Cresselia \texttt{sv6pt5}-21 and Team Rocket's Bother-Bot \texttt{sv10}-172), extra turns,
markers, and the Lost Zone. The universal construction's key ingredient, a rest-of-game modifier
granted again and again, therefore does not occur in Standard, and neither does a once-per-game
attack that grants one.

Other mechanisms are present. Thirty-four cards carry a damage reduction into the opponent's
turn, for example Registeel ex (\texttt{me2pt5}-145), ``During your opponent's next turn, this
\pkmn{} takes 50 less damage from attacks (after applying Weakness and Resistance).'', and the
Item Iron Defender (\texttt{me1}-118, \texttt{me5}-110); five attacks heal by the damage done
(Virizion \texttt{rsv10pt5}-10, Trevenant \texttt{sv6}-13, Iron Moth \texttt{sv6pt5}-9,
Volcarona \texttt{sv8}-25, Paras \texttt{sv9}-4). So Standard decks are in general not tame,
and the faithful theorem is the one that bears on them. Seventy cards return cards from the
discard pile (65 found by regular expression, 5 more by reading), and four \pkmn{} can attack
twice, all through the Ability Festival Lead with the Stadium Festival Grounds in play (Dipplin
\texttt{sv6}-18, Goldeen \texttt{sv6}-44, Swirlix \texttt{sv6}-89, Seaking
\texttt{sv8pt5}-21: ``If Festival Grounds is in play, this \pkmn{} may use an attack it has
twice.''). These are two of the three absences on which the faithful bound rests; the third,
extra turns, does not occur. The faithful bound therefore does not yet cover Standard. Also
present are Abilities on 316 \pkmn{}, seven of them usable ``as often as you like'' (for example
Emboar \texttt{rsv10pt5}-13, which attaches Basic Fire Energy from the hand); 24 attacks that
flip a coin until tails; and HP that grows with attached Energy (Conkeldurr
\texttt{zsv10pt5}-49, ``This \pkmn{} gets +40 HP for each Fighting Energy attached to it.'').
Every Standard \pkmn{} except 58 Dragons has Weakness $\times 2$; Resistance is $-30$ on 401
cards and $-20$ on one.

Two findings bound what the model over-approximates. No Standard card raises an existing Poison
as Croagunk does, and the condition payloads (up to 16 damage counters at each Checkup, Mega
Dragalge ex \texttt{me4}-65) are replaced, not added, by a later Poison. And as far as the
search's patterns reach, no Standard effect replays a same-turn damage bonus within a turn: the
bonus Supporters are limited to one per turn, and the one bonus Item, Premium Power Pro
(\texttt{me1}-124), is returned to the hand within a turn by no Standard effect except an
attack (Dedenne \texttt{sv8}-87), which ends the turn; this is an argument from the search, not
a theorem.

\subsection{Where memory could hide}\label{sec:hide}

The survey rules out an unbounded store that card text reads beyond a yes-or-no test, and
Theorems~\ref{thm:finite} and~\ref{thm:xfinite} rule one out in the model for tame and for
faithful decks. Here we examine the places
in the printed game where a store is most plausible. Each claim carries its evidence:
\ev{regex} for a search of the card data by regular expression, \ev{hand} for an argument by
hand in this paper, and \ev{Lean} for a theorem of the model.

\paragraph{Damage counters.}\label{sec:hide-stack} A pile of damage counters on one \pkmn{} is
the natural candidate for a register. The card pool supplies an increment (attacks and
effects that place damage counters), a decrement (healing, and Alakazam's Damage Swap, which
moves damage counters one at a time ``as often as you like during your turn''), 22 wordings of a
zero test (for example ``full HP'' or ``no damage counters''), and exact tests such as Mega Absol
ex's Terminal Period (``If your opponent's Active Pok\'emon has exactly 6 damage counters on
it, that Pok\'emon is Knocked Out'') \ev{regex}. One card, a Lunala (\texttt{cel25}-15), doubles
the damage counters; we found no effect that halves them or tests their parity. Damage
therefore has the interface of a Minsky register. What it lacks is capacity: the number of
damage counters is below the \pkmn{}'s HP at every Knock Out check \ev{hand, from the rules of
Section~\ref{sec:game}}. Printed HP is at most 380, so an ordinary \pkmn{} holds at most 37
damage counters; the store search estimates about 250 for a Conkeldurr whose HP grows with
attached Energy, which is still bounded by the 60 cards of a deck. In the model, every move
other than a Knock Out starts with every \pkmn{} below its HP, and a \pkmn{} at or above its HP
is untouched until it is Knocked Out \ev{Lean: Theorem~\ref{thm:untouched}}.

\paragraph{The single attack.}\label{sec:hide-attack} Outside an attack or a Checkup, a
\pkmn{} whose damage reaches its HP is Knocked Out at once (the Advanced Player's Rulebook and
rulings 85 and 449 \cite{apr,compendium}), and Checkup applies a fixed set of updates once and
then checks for Knock Outs. The resolution of one attack is therefore the only window in which
damage may exceed HP while a player's choices are carried out \ev{hand}. Computing in that
window needs an attack that loops over a state it reads. Copy attacks (21 wordings, such as
``Choose 1 of the Defending Pok\'emon's attacks and use it as this attack'') each select one
attack and resolve it in place of themselves, and no text uses another attack afterwards or
repeats one \ev{regex}; a copy attack can copy another copy attack, and the control-flow search
found that each level adds a choice and that the depth is bounded by the \pkmn{} in play.
Flip-until-tails loops count heads and change no state in between; the one variant that reads
the state, Mamoswine's Ramming Strike, restarts at a cost in its own damage and is bounded by
its own HP \ev{regex, census by eye}. The ``as often as you like'' sentences that the census
lists, such as Damage Swap, are used outside an attack, but we did not check all 83. As far as
the search reaches, then, one printed attack is straight-line code with nondeterministic
choices \ev{hand}; we found no ruling on two copy attacks that could copy each other. The model
is more permissive, since an attack's effect may put \lean{ite} inside \lean{iter}, and still,
for tame decks, whatever one attack computes must fit in the clamped state \ev{Lean:
Theorems~\ref{thm:notsim}, \ref{thm:clamp} and~\ref{thm:finite}}.

\paragraph{Special Conditions.}\label{sec:hide-cond} The Special Conditions give each Active
\pkmn{} $4 \times 2 \times 2 = 16$ combinations, and only the two Active \pkmn{} carry them
\ev{rules audit; Lean: \lean{rot}, \lean{poisoned}, \lean{burned}, and \lean{WellFormed}}.
Poisoned and Confused persist until a cure, a retreat or an evolution; Paralyzed, Asleep and
Burned are timers that Checkup may end. They are finite control, not memory, and in the model
they are among the flag values enumerated by the list of Theorem~\ref{thm:finite} \ev{Lean}.

\paragraph{Poison strength.}\label{sec:hide-poison} Croagunk's Poison Up (``If the Defending
Pok\'emon is Poisoned, put 3 more damage counters on that Pok\'emon between turns. This effect
can be applied more than once.'') makes Poison strength a number that grows by one per use
\ev{regex, census by eye}. As a register it is weak in three ways. It only increments, and it
resets to zero when Poison ends or a new Poison replaces it \ev{store search; rulebook p.~16;
Lean: \lean{cure}, \lean{retreat}, \lean{evolve} and a new Poison set \lean{poisonUp := 0}}. Only Checkup reads it \ev{store search; Lean:
\lean{Test} has no case for it}; the survey records no card that tests it, but no part of the
survey searched for such a test. And Checkup reads it destructively: the damage it adds meets
the Knock Out threshold, so all values large enough to Knock the \pkmn{} Out are
indistinguishable, and reading one discards the \pkmn{} and the number with it. The store
search argues that in the printed game Poison strength rises at most once per turn and is read
at the next Checkup, so the \pkmn{} is Knocked Out while it is below about a third of its HP in
damage counters \ev{hand}; we have not checked this against effects that skip Checkup, such as
the extra turns of Medicham~V's Yoga Loop. In the model, where \lean{poisonUp} can be raised
without bound inside one effect, the excess above HP is unobservable \ev{Lean:
Theorem~\ref{thm:clamp}}.

\paragraph{Control flow.}\label{sec:control} Control is cheap; memory is the obstacle. The
control-flow search found that every main-phase action is optional, including attacking, and
that no card forces a player to attack. The mandatory surface is the draw at the start of each
turn, Checkup, triggers on a player's optional action, and a handful of state triggers (Neo
Genesis Berries, Star Piece, Team Aqua and Magma Belts, Brock's Primeape's ``Scram'', Ditto's
``Transform''); Checkup is the only loop body that runs by itself, and no configuration was
found in which the game runs indefinitely with no decisions. The pool nonetheless offers
threshold and equality tests on damage, HP, hand size and Prizes, parity tests (Mr.\ Mime-GX's
Magic Odds), a multi-way branch on hand size (Kricketune), timers (Paralysis), derandomizers
that turn coin flips into tails (Shiftry, Malamar), and locks on Items, Abilities, evolution and
retreat. When one player executes a program, that player is the program counter, so the
realistic assumptions about play are cooperative play or a unique legal line enforced by locks.
Two clocks end a game: a player who cannot draw at the start of their turn loses
\cite{rb-ssp}, and taking the last Prize card wins. An unbounded game must return cards to the
deck at least once per turn per player, and almost every Knock Out yields a Prize; only a few
\pkmn{} (Shedinja's Empty Shell, Robo Substitute, Lillie's Pok\'e Doll) give none. These clocks
bound the time of particular strategies, not the memory: in the model, \lean{handToBottom}
already lets a player refill the deck, games of unbounded length exist, and
Theorem~\ref{thm:finite} holds regardless. A game of unbounded length over a finite state space
is a finite automaton running forever.

\subsection{What remains open}\label{sec:linger-next}

The printed decks are not all tame, since they can carry a reduction (Iron Defender), and in
today's Standard pool reductions and \lean{drain} are common (Section~\ref{sec:standard}). They
do not supply the universal construction, for the reasons that follow. Their modifiers fall, as
far as the model expresses them, under the faithful theorem (Theorem~\ref{thm:xfinite}), whose
limits Section~\ref{sec:faithful} lists. This subsection states what that leaves open; nothing
in it is proved.

\paragraph{Why the construction does not run on printed cards.}\label{sec:linger-printed} The
targeted reading read every sentence with a rest-of-game duration and every sentence with
``your next turn'' or ``your opponent's next turn'', and classified 1,616 distinct sentences
with a duration cue by regular expression. Four templates last for the rest of the game. Three
come from GX attacks (Lucario \& Melmetal-GX, 30 less damage to the player's Metal \pkmn{};
Arceus \& Dialga \& Palkia-GX, 30 more damage; Latios-GX, no GX attacks for the opponent), and
the fourth is a Trainer that turns Prize cards face up. A player may use only one GX attack in
a game, so each player's rest-of-game damage modifier takes at most one fixed value, and the
construction, which grants a rest-of-game modifier in every round and grants reductions to the
opponent, cannot run on these cards. Today's Standard pool has none of the four, and no other
rest-of-game modifier. Bonuses and reductions do accumulate within a turn: Items
are unlimited per turn (PlusPower $+10$, Electropower $+30$, Iron Defender $-30$), and
once-per-turn Abilities and static effects add more. The targeted reading bounds each such
total by the copies of each card (at most four per name) times a printed constant, found no
printed effect that accumulates across turns without a bound by cards, and found none of the
143 ``as often as you like'' and ``any number of times'' sentences producing a lingering damage
or HP modifier. This bound is an empirical claim about the card pool, and it holds only if no
card-recycling loop replays a modifier card within one turn. The rules audit records, from
secondary sources not checked against card text, one such loop in the Expanded format, the
``Infinite Flapple Loop'', in which Flapple's Apple Drop shuffles Flapple back into the deck and
other cards let the player cycle the deck without limit in one turn; Flapple was then banned
from Expanded \cite{flapple}, in an announcement that does not state the reason. We have not
checked whether such a loop can replay a ``+$N$ damage this turn'' card. If one can, the
in-turn bonus is unbounded, though it still resets at the end of the turn.

\paragraph{A second route to universality (argued).} The three slots behave differently at a
turn boundary: this turn's modifiers expire, carried modifiers are written in one turn and read
in a later one, and rest-of-game modifiers never expire. One might expect that if an attack,
which happens once per turn, is the only reader of the difference between bonus and reduction,
then only boundedly many bits cross each turn boundary. That is false as stated, by a
guess-and-verify construction that needs no rest-of-game slot and one exact test per turn. Hold
a register $c$ as a carried reduction. The player to move guesses a number $m$ and plays, $m$
times, a replayable effect that adds $1$ to this turn's bonus and $k$ to a carried slot; the
attack then tests that the damage done is exactly a fixed $n$, which forces $m = c$, and the
carried slot then holds $k \cdot c$. Division with remainder is similar, and multiplication and
division by constants with a zero test suffice for a universal machine with one register
\cite{minsky}. The construction needs the two writes to be \emph{coupled}: done by one
replayable effect, so that their counts cannot drift apart. This is a hand argument that no
second reader has checked. The one replayable carried damage modifier we found on printed cards
is not coupled: a regular-expression search (Trainer cards, and Abilities without a
once-per-turn limit, for modifiers whose duration outlasts the turn) found only the Item Iron
Defender (me1-118, me5-110), ``During your opponent's next turn, all of your Metal Pok\'emon
take 30 less damage from attacks from your opponent's Pok\'emon'', which writes a carried
modifier only. On this evidence a carried reduction is bounded by four copies times 30 unless a
recycling loop replays the Item, in which case it is unbounded but not coupled. We have not
checked whether two copies stack, whether a loop can couple two different cards, or attacks
reached through extra turns. In the model the route is closed: inside one effect \lean{LoopFree}
forbids it, and across effects the replays it needs do not exist, since an attack writes a
bounded amount once per turn and an Item or Supporter is played at most once per game. Card
recycling or extra attacks would reopen it.

\paragraph{The open question.} Whether finiteness survives once card recycling, repeated attacks
(``attack twice'', ``attack again'') and extra turns are in the model is open. Each exists on
printed cards, and each reopens the store bound of Theorem~\ref{thm:xfinite}. For today's
Standard pool the question is concrete: its decks are not tame, so the faithful theorem is the
relevant one, and Standard has card recycling (70 cards) and attacking twice (the four Festival
Lead cards) but no extra turns. The open case is therefore the faithful theorem with card
recycling and repeated attacks in the model, for the Standard pool. A proof would need a bound
on the replays these effects allow, or a coarser quotient than the exact-store clamp; the
argument of Section~\ref{sec:standard} that no Standard effect replays a same-turn bonus within a
turn is a candidate for the first.

\subsection{Conjecture and relaxations}\label{sec:conjecture}

\begin{conjecture}\label{conj:main}
Take the English card pool, without Japanese-only cards and without \pkmn{} made of several
cards, under the game's current rules.
\begin{enumerate}
\item[(i)] For every pair of legal decks drawn from the pool in which no card-recycling loop can
  replay a card within one turn, cap damage and Poison strength on each \pkmn{} at its current
  HP, and every damage bonus and reduction at a constant computed from the decks. Then the
  capped states reachable from a set-up position lie in a finite set, and two reachable states
  with the same capped image have the same moves up to that image.
\item[(ii)] For every pair of legal decks drawn from the pool, the game is not Turing complete
  under any assumptions about play whose memory is the game state.
\end{enumerate}
\end{conjecture}
\noindent For the decks of (i), (ii) follows from (i), since the capped game is then a finite
transition system with the same moves. For decks with a recycling loop, (i) can fail within a
turn while (ii) still holds. A proof of the conjecture needs three things beyond this paper: a
model that covers the current rules, an effect language that expresses every printed card text,
and an answer to the open question of Section~\ref{sec:linger-next}, a finiteness theorem in the
manner of Theorem~\ref{thm:xfinite} for a model with card recycling, repeated attacks and extra
turns, with its hypothesis discharged for the card pool (Appendix~\ref{app:future}). Today's
Standard pool is the nearest case: it needs card recycling and repeated attacks but not extra
turns, and it has neither rest-of-game modifiers nor a Croagunk-style raise of Poison strength,
the printed counterparts of the generalisations behind Theorems~\ref{thm:universal}
and~\ref{thm:infinite}.

\paragraph{Relaxations.}\label{sec:lattice} A finiteness conjecture invites the converse
question: which change to the rules would make the game universal? We record candidates,
strongest evidence first. Only the first has a machine-checked part.
\begin{itemize}
\item \emph{Rest-of-game modifiers in attacks that are not GX attacks, with the generality of
  Section~\ref{sec:counter-generality} (machine-checked for the model).} Under these
  departures the model simulates every two-counter machine (Theorems~\ref{thm:universal}
  and~\ref{thm:sharp}). In the printed game the change would need new cards or a lifted
  once-per-game rule.
\item \emph{A coupled replayable writer (argued).} A replayable effect that writes a this-turn
  and a carried modifier together would give universality without the rest-of-game slot, given
  card recycling or repeated attacks to replay it (Section~\ref{sec:linger-next}).
\item \emph{HP as a parameter (conjectured).} Printed card text already gives damage counters
  an increment, a decrement and a zero test (Section~\ref{sec:hide}); only capacity is missing.
  We conjecture that with unbounded HP two damaged \pkmn{} serve as the two registers of a
  Minsky machine under cooperative play, with the Special Conditions and the configuration of the
  board as finite control. With HP a finite parameter $n$, the same construction would simulate
  two-counter machines whose registers stay below $n$. A proof would have to check that the
  increment, decrement and zero test can be aimed at the two chosen \pkmn{}, that the zero test
  can steer control rather than merely trigger a Knock Out, that one instruction fits the
  per-turn limits, that both clocks of Section~\ref{sec:control} can be defeated indefinitely,
  and that every card used is legal together under one set of rules. This item has no checked
  part; it needs no new kind of card, only larger HP.
\item \emph{Deck size as a parameter (conjectured).} With decks of unbounded size the deck is a
  queue: Caitlin, Kofu and Claydol's Cosmic Power put cards from hand on the bottom of the deck
  in a chosen order, and the draw removes cards from the top. A cooperative player could then
  run a tag system, with deck-out as halting, provided the construction avoids the many effects
  that shuffle the deck.
\item \emph{A count reader for markers.} A card that reads the number of Imprison markers,
  together with one that removes them one at a time, would make the marker tally a register and
  would remove the justification of Lemma~\ref{thm:flag}.
\item \emph{An unbounded Bench}, which might hold a row of cells.
\end{itemize}

%======================================================================
\section{Related work}\label{sec:related}
%======================================================================

\paragraph{Card games.} Magic: The Gathering is Turing complete under forced play, using
creature tokens as tape \cite{cbh-mtg}; a follow-up shows that optimal play is
non-arithmetic \cite{biderman-arith}, and a later paper, which we have not read, embeds a
programming language in Magic \cite{yin-churchill}. Yu-Gi-Oh!\ \cite{ygo} and Riftbound
\cite{riftbound} are discussed in Section~\ref{sec:intro}. A self-published argument claims
that Netrunner is Turing complete under ``choose or lose'' constraints \cite{netrunner-tc}; we
have not verified its content. Lower bounds without universality include weak NP-hardness of
Netrunner mate-in-1 or -2 \cite{bosboom-netrunner}, PSPACE-hardness of perfect-information
Hearthstone \cite{zhang-hearthstone}, and NP-hardness of single-player UNO \cite{uno}. Turing
completeness has also been shown for Sid Meier's Civilization \cite{dewynter-civ}, a video
game. Results on the \pkmn{} \emph{video} games \cite{nintendo-hard} concern a different game.
Our search found no negative argument (``finite state, so not universal'') for the \pkmn{} TCG
in any source; we did not search for negative results about other card games. A negative
result is only as good as the model it is proved for, which is why we mechanize the model and
list its departures.

\paragraph{The \pkmn{} TCG in Lean.} Ramos and Soria \cite{ramos-soria} describe a Lean~4
development of the \pkmn{} TCG aimed at metagame analysis: a matchup matrix over 14
archetypes, Nash equilibria and replicator dynamics. By their account the artifact has 87
files, about 31,900 lines and 2,627 theorems, with no \lean{sorry} or custom axioms, and 145
theorems use \lean{native\_decide} in the game-theory layer. As described in their
Section~III-B, the rules layer includes a game state with per-player zones, turn ownership and
a finite phase machine, a deck-legality checker, card-count conservation, Bench-size bounds,
Prize accounting, one card effect (\lean{professorsResearchEffect}) with a per-player card
count theorem, and a type-effectiveness table. In the part of their text we read, we found no
step relation, no Knock Out or damage semantics, no effect language, and no statement about
finiteness, decidability, Turing completeness or complexity. We read roughly 80,000 of the
106,000 characters of the arXiv HTML verbatim and the rest only as a summary, and we did not
parse the PDF. We have not inspected their Lean source: the paper calls it anonymized
supplementary material and cites an IEEE DataPort record whose files we found behind a
subscription login, so we cannot compare definitions. Their work is prior art for Lean and the
\pkmn{} TCG. Measured against their paper's description, and subject to what their source may
contain, what this paper adds is an operational semantics with an effect language, and the
two finiteness theorems and the reduction proved over it.

%======================================================================
\section{Conclusion}\label{sec:conclusion}
%======================================================================

\paragraph{Finiteness for tame decks.} For a Lean model of the \pkmn{} TCG with the rulebook's
turn structure, typed costs, Weakness and Resistance and damage calculation, Special Conditions,
Checkup, Knock Out timing, markers, lingering damage modifiers and effects that read damage,
tame decks from a well-formed start with no reduction in effect have a finite clamped game with
the same moves as the game itself; the clamp caps stored bonuses at twice a number above every
printed HP, Resistance and test constant. Effects can observe damage above HP inside one attack,
but no move carries it further.

\paragraph{Finiteness for faithful decks.} A clamp that keeps every modifier exact is a two-way
bisimulation for every deck. For every well-formed start whose cards put no damage bonus or
reduction under \lean{iter} and write rest-of-game bonuses or reductions in attacks only through
once-per-game GX attacks, the reachable states fall into finitely many of its classes; damage
reduction, \lean{drain}, carried modifiers and Trainer-card modifiers are all allowed. The bound
rests on the model's lack of card recycling, repeated attacks and extra turns.

\paragraph{The universal side.} Dropping the GX restriction admits decks, which no printed
card supplies, that simulate every two-counter machine: from a well-formed start a halt marker
is reachable by some sequence of legal moves if and only if the machine halts, and soundness
holds over all play. With Minsky's theorem, which we cite and do not formalise, reachability of
the halt marker over these starts is undecidable.

\paragraph{What remains open.} Whether finiteness survives once card recycling, repeated
attacks and extra turns are in the model is open. For today's Standard pool, whose decks are
not tame and which has card recycling and attacking twice but no extra turns, the open case is
the faithful theorem with recycling and repeated attacks modelled. A second route to
universality, an unchecked hand argument, would need a replayable card writing two coupled
modifiers and the replays that those effects supply.

\paragraph{The printed card pool.} A survey of 20,776 English card records found no unbounded
store that card text reads beyond a yes-or-no test, with one case open, a card-recycling loop
that replays a damage modifier within one turn; the likeliest hiding places, damage counters,
the single attack, Special Conditions and Poison strength, are each bounded or read only
destructively. Today's Standard pool has no rest-of-game modifiers, so the universal
construction has no Standard instance.

\paragraph{Conjecture.} We conjecture that the printed \pkmn{} TCG is not Turing complete.

\paragraph{Availability.} The Lean development is public \cite{lean-pokemon}; everything cited
here is on its branch \texttt{master} at commit \texttt{ce0de3b}, and the two-counter
construction is in \texttt{Pokemon/Counter.lean}. The survey files and census scripts are not
yet published.
% TODO: decide where the survey, census scripts and template table will be published, and cite them.

\paragraph{Acknowledgement.} The Lean development was reviewed adversarially by independent
automated reviewers during development.

\bibliographystyle{plain}
\bibliography{refs}

\appendix

%======================================================================
\section{The departures from the rulebook}\label{app:departures}
%======================================================================

The header of \texttt{Rules.lean} gives the sources of the model's turn and damage rules and
lists every known place where the model is simpler than the rulebook. We restate both,
condensing the wording; Table~\ref{tab:departures} groups the list.

\paragraph{Sources of the turn and damage rules.} The rulebook quoted is the 2023 edition
\cite{rb-151}; the rules audit read the 2024 rulebook \cite{rb-ssp}, and the 2026 rule changes
were read through a summary \cite{pp26}. The model follows the latest source, and for the rules
below the three agree except where noted. An attack ends the turn (pp.~9, 13). A failed
Confused attempt also ends the turn: the rulebook says only that ``the attack doesn't happen''
(p.~16), and the model reads the Confusion check as step~C of using the attack (p.~20). The
first-turn bans are the rulebook's (pp.~11--13): no attack and no Supporter for the starting
player on the first turn of the game, no evolving on either player's own first turn. The damage
calculation stops at base damage 0 (p.~20), once per attack, applied to the base damage as the
attack's text computes it. A new Poison replaces the old one (p.~16); that Croagunk's raised
Poison strength ends with the replaced Poison is the model's reading, and we found no ruling on
it. There are nine Basic Energy types (glossary, p.~39) and one Stadium in play at a time
(p.~12). For simultaneous wins the 2023 rulebook has Sudden Death (p.~21), which the 2026
changes replaced by tiebreaker games, unverified against a rulebook; the model has neither.
Knock Outs at times other than after an attack or Checkup appear in no English rulebook; the
model's rule comes from the Advanced Player's Rulebook \cite{apr}. Typed costs follow the
rulebook's example of Quaxwell's Spiral Kick (p.~13): a typed symbol needs one Energy of that
type, while for a Colorless symbol ``any type of Energy can be used''. Weakness and Resistance apply to the Active \pkmn{} only (p.~14), and the damage
calculation is the rulebook's six steps (p.~20). Item and Supporter effects do no attack damage,
since ``damage counters aren't affected by Weakness, Resistance, or any other effects'' (p.~20,
step~1) and printed Trainer cards place damage counters; a printed Trainer card whose text says
``does damage'', if there is one, would lose any modifier the rulebook applies to it, and no
search for such cards was made.

\paragraph{Known gaps.}
\begin{enumerate}
\item Setup, mulligans and the opening coin flip are not modelled; play starts from a set-up
  state.
\item Energy is attached to, and evolution applied to, the Active \pkmn{} only.
\item Attack costs are typed: each typed symbol is paid by a unit of attached Energy of that
  type and each Colorless symbol by any unit, every unit paying one symbol (\lean{Pays}). Only
  Basic Energy exists, one unit of its type. Special Energy (several units, or units of several
  types), cost changes (an attack that costs one Colorless Energy less) and Energy that provides nothing
  in some states are absent. A Retreat Cost is paid by discarding that many Energy cards, which
  with Basic Energy only is the rulebook's rule.
\item Weakness is $\times 2$ to one type and Resistance $-N$ to one type, as printed today
  (every Standard Weakness is $\times 2$; Resistance is $-30$, once $-20$). Absent are the
  ``$+10$'' to ``$+40$'' Weakness of older cards, two Weaknesses or two Resistances on one card
  (a \pkmn{} has one type here), and effects that change or ignore Weakness or Resistance.
  Weakness and Resistance never apply to the Bench, which attack damage does not reach in the
  model.
\item Knock Outs are separate moves, forced by \lean{Settled} before any other move. They are
  taken one at a time, in either order when several are pending, so when both players take
  their last Prize from the same attack or Checkup the first Knock Out taken decides the game
  (no Sudden Death, which the 2023 rulebook has, nor the tiebreaker game that replaced it in
  2026). A player with no \pkmn{} to promote loses at once. Only Active \pkmn{} can be Knocked
  Out, which suffices while no effect damages the Bench. A Knock Out always gives exactly one
  Prize.
\item The source of the Knock Out timing outside attacks and Checkup: we found no rule in the
  2023 English rulebook that Knock Outs are checked at once at other times, and the 2024
  rulebook as the rules audit quotes it has none either. The model's immediate Knock Outs in
  mid-effect (\lean{TRun}) and its \lean{Settled} premise on every other move rest on the
  Advanced Player's Rulebook of \pkmn{} Asia \cite{apr}, a document of another region's rules
  text. It is a modelling decision until an English source confirms it.
\item Knock Outs in mid-effect (Item and Supporter cards, \lean{TRun}) are checked after each
  primitive step. A printed sentence that the model writes as several primitives therefore gets
  a check between them, and the outcome in which the whole sentence acts on the original
  \pkmn{} is absent: \lean{seq (counters 5) (scale opp 1)} against a 5-HP Active always acts on
  the promoted \pkmn{} in its second step. The replacement is promoted at once, before the next
  step; no ruling was found on when it is chosen in mid-effect. The choice of \pkmn{} is
  nondeterministic, which covers the owner's choice, but outcomes in which later steps find an
  empty Active Spot are absent. When a Prize ends the game, or its owner has no \pkmn{} to
  promote, the rest of the effect is not resolved. Simultaneous Knock Outs cannot arise inside a
  Trainer effect, because each primitive step damages at most one Active \pkmn{} and starts
  settled; a printed effect that damages several \pkmn{} at once is not expressible.
\item A Knock Out takes the top card of the Prize list, whereas the rulebook lets the player
  choose any Prize card (a finite choice).
\item HP is the printed value: there are no HP modifiers.
\item Markers are flags, and the Coloring marker is a type. The printed Imprison and Shock-wave
  counts are uncapped, but only ``has any'' is read and removal takes them all, so a flag is
  exact for them (Lemma~\ref{thm:flag}); this is a modelling assumption about the card pool.
  Effects may place a marker or a Coloring marker on either player's \pkmn{}, whereas printed
  cards place them on the opponent's, and \lean{unmark} may clear either player's, whereas
  Tropius clears its own player's; marker effects are general primitives, usable inside
  \lean{iter}. Every marker stays through evolution, retreat and promotion and leaves play with
  its \pkmn{}; card text settles this only for Char. Checkup reaches only the Active \pkmn{},
  so a Char marker on a Benched \pkmn{} does nothing until it is Active again, whereas the
  printed Char also damages a Benched \pkmn{}. Absent are the readers of Imprison (they disable
  Pok\'e-Powers and Pok\'e-Bodies, and the model has no Abilities), tests on Benched \pkmn{}'s
  markers, Dark Ivysaur's transient markers allocated by the opponent, and the marker-like states
  that the targeted reading lists (Poison payloads, delayed triggers, attack locks bound to one
  \pkmn{}, history flags, face-up Prizes).
\item Effects read only the two Active \pkmn{}: their damage (\lean{Test.atLeast},
  \lean{Test.exactly}, \lean{scale}), their Special Conditions (\lean{Test.cond}), their
  markers (\lean{Test.marked}) and their type (\lean{Test.isType}). There are no tests or
  counts on Benched \pkmn{}, Energy, cards in hand, deck, discard pile or Prizes, Poison
  strength, or HP (``if it has full HP'', ``remaining HP''), or on the lingering modifiers in
  effect. \lean{scale} writes only to the opponent's Active \pkmn{} (Lunala doubles the damage
  on each of the opponent's \pkmn{}, Bench included). The damage step's amount reads the damage
  on either Active \pkmn{} (\lean{Amt.per}), and the step itself reads the Weakness and
  Resistance of the opponent's Active \pkmn{} and the type of the player's.
\item Lingering modifiers (Section~\ref{sec:linger-mech}): amounts add up without limit, since
  \lean{linger} is a general primitive usable any number of times and inside \lean{iter}.
  Printed same-turn bonuses and reductions are bounded by copies times a constant if no
  card-recycling loop replays them, an empirical claim about the card pool; rest-of-game ones
  come from once-per-game GX attacks; a reduction carried into the opponent's turn exists on a
  printed Item, Iron Defender. Modifiers belong to a player, not to one \pkmn{}, so ``this
  \pkmn{} takes 30 less damage during your opponent's next turn'' protects whichever \pkmn{}
  that player has Active, and does not end when the \pkmn{} leaves the Active Spot. Damage done
  by an attack is computed once per attack, by the rulebook's steps (\lean{dealtBy}): nothing if
  the base damage $n$ is 0; otherwise $n$ plus the attacker's bonus, doubled under Weakness, less
  Resistance, less the defender's reduction, each subtraction stopping at 0, or nothing under
  prevention; \lean{counters} is not modified, as in the rulebook. The bonus enters at step~2
  and the reduction at step~5, as printed modifiers say; modifiers that name the other step are
  not expressible. Threshold prevention, prevention of effects, locks on Supporters, Stadiums,
  Abilities or Energy attachment, cost changes, effects that change Weakness, Resistance or type
  (other than the Coloring marker), ``attack twice'', ``take another turn'' and ``that much more
  damage'' are absent. Finiteness is proved for tame decks and for
  faithful decks (Section~\ref{sec:finite-side}); for decks that no printed card supplies, one
  construction, which uses rest-of-game modifiers in attacks that are not GX attacks together
  with the generality of Section~\ref{sec:counter-generality}, simulates every two-counter
  machine (Section~\ref{sec:linger}).
\item An attack's damage: an attack with a printed base damage and no damage text does that
  damage and then its effect (\lean{Attack.script}); an attack whose text computes or conditions
  its damage has base damage field 0 and one damage step in its effect, at most one on each line
  of resolution (\lean{Effect.oneHit}), whose amount is evaluated once when the step is reached
  (Jynx, \texttt{base1}-31, is \lean{damage (add (const 2) (per opp 1))}). Effects before the
  damage step are the rulebook's ``effects that happen before damage''; effects after it see the
  damage placed. The amount language has no tests of its own: a damage that depends on a test, a
  coin flip or a choice is written with the damage step inside the branches of \lean{ite} or
  \lean{choice}. Not expressible: damage counted from anything other than damage counters on an
  Active \pkmn{} (Energy, Benched \pkmn{}, cards in hand or discard pile), and two damage
  calculations in one attack. The generality given up: an attack cannot make several damage
  steps (an \lean{iter} of \lean{damage}, which would apply the modifiers per step and is not the
  rulebook's calculation), and Item and Supporter effects cannot do attack damage (their printed
  counterparts place damage counters).
\item \lean{ite}, \lean{scale}, \lean{selfDamage} and \lean{healAll} are general, usable in
  any effect, any number of times and inside \lean{iter}, whereas printed cards use them once,
  in fixed forms. The model over-approximates the printed texts.
\item GX attacks (\lean{Attack.gx}): a player who has used one (\lean{gxUsed}) cannot use
  another in the same game, and \lean{Move.attack} sets the flag. A Confused \pkmn{} can try a GX
  attack only under the same condition, and a failed attempt does not use it up; this answers an
  unverified ruling. VSTAR powers are absent.
\item Special Conditions: effects apply them only to the opponent's Active \pkmn{}
  (\lean{inflict}) and remove them only from the player's own Active (\lean{cure}). The
  variants with a number (``put 2 damage counters instead of 1 on the Poisoned \pkmn{}'') are
  absent. A new Poison replaces the old one, as the rulebook says, so poisoning a \pkmn{} that
  is already Poisoned resets its Poison strength to 0 (\lean{InPlay.inflict}). Croagunk's
  cumulative increase is modelled as a number on the \pkmn{} (\lean{poisonUp}) that belongs to
  the Poisoned Condition: it lasts until Poison ends or is replaced, and stacks across uses of
  the attack, which raises an existing Poison and does not re-Poison. A Confused \pkmn{} whose
  flip is tails ends the turn as an attack does, and its attack must be one it could otherwise
  use.
\item \pkmn{} Checkup covers the five Special Conditions and the Char marker only: no Ability,
  Stadium or other ``between turns'' card effect exists in the model, so the rulebook's choice
  of order between Special Conditions and those effects does not arise. Burn uses the current
  rule (2 damage counters, then heads cures). Paralysis ends at the Checkup after its owner's
  turn, which is exact here because only the opponent's \pkmn{} can be Paralyzed. Coin flips
  are nondeterministic choices, so the model is silent about probability.
\item \lean{poisonUp n} is a general primitive, usable in any effect and inside \lean{iter},
  whereas the printed card raises Poison strength once per use of its attack. The model
  over-approximates it, and one Item can raise Poison strength without bound
  (Theorem~\ref{thm:infinite}).
\item Effects damage only the two Active \pkmn{}: \lean{counters} and \lean{scale} damage the
  opponent's, and \lean{selfDamage} the player's own. No effect damages the Bench.
\item \lean{heal} and \lean{healAll} remove damage from the player's own Active \pkmn{} only,
  not from the Bench nor from the opponent's \pkmn{}.
\item Stadium and Tool cards occupy state but have no move. The Stadium slot is one slot shared
  by both players, with the card's owner, as the rulebook has it, but it only holds a card of the
  start state, so the rules that come with playing a Stadium (one per turn, no Stadium of the
  same name, the old one discarded) and Tools (one per \pkmn{}) have nothing to act on.
  Abilities are absent.
\item Start states are not required to be well formed. From a \lean{WellFormed} start no
  reachable state without a winner is stuck (Theorem~\ref{thm:progress}), but from other starts
  (for example, no winner and an empty Prize list) a pending Knock Out may be impossible to
  take. Nor are start states required to have consistent turn flags: from a start in which it
  is the first turn of the game and the player to move has already had a turn, that player
  could evolve on the game's first turn. Consistency is \lean{GameState.FlagsOk}; play
  preserves it, and the first-turn bans are theorems under it (Theorem~\ref{thm:turn}).
\item The first-turn rules are the rulebook's: the starting player cannot attack or play a
  Supporter on the first turn of the game, and neither player can evolve on their own first
  turn. Card text that lifts these bans (``If you go first, you can use this attack during your
  first turn'', Delibird \texttt{me1}-105; 26 Standard cards) is not expressible.
\end{enumerate}

%======================================================================
\section{State inventory of the printed game}\label{app:inventory}
%======================================================================

Table~\ref{tab:inventory} condenses the rules audit's inventory of the state of the printed
game. The labels give the source of each bound: \textsf{Q} is quoted rules text, \textsf{R} an
official ruling, \textsf{I} an inference of ours, and \textsf{U} an unverified or secondary
source.

\begin{table}[ht]
\centering
\small
\begin{tabularx}{\textwidth}{@{}l >{\raggedright\arraybackslash}p{2.6cm} >{\raggedright\arraybackslash}X@{}}
\toprule
Component & Bounded? & Bound and source \\
\midrule
Deck (ordered) & yes & $\le 60$ cards; order is a permutation (\textsf{Q} \cite{rb-ssp}) \\
Hand & yes & no hand-size limit; bounded by card conservation (\textsf{I}) \\
Discard pile, Lost Zone & yes & cards of the fixed multiset (\textsf{I}) \\
Prize cards & yes & 6 at setup; effects move cards in and out (\textsf{Q} \cite{rb-ssp,apr}) \\
Active, Bench & yes & one Active; Bench $\le 5$, raised to at most 8 by a few cards (\textsf{Q}/\textsf{U}; the survey found no larger Bench) \\
Stadium & yes & one in play (\textsf{Q} \cite{rb-ssp}) \\
Attached cards & yes & no limit on Energy per \pkmn{} (\textsf{Q} \cite{apr}), but they are cards \\
Damage counters & yes, via HP & Knock Out at damage $\ge$ HP; overshoot is transient (\textsf{Q}/\textsf{R}, Section~\ref{sec:game}) \\
Current HP & yes, if every modifier is bounded & modifiers found are tied to cards in play (\textsf{Q} \cite{apr}); see Section~\ref{sec:findings} \\
Special Conditions & yes & five; Burned and Poisoned replace rather than stack (\textsf{Q} \cite{rb-ssp,apr}) \\
Once-per-game, once-per-turn flags & yes & finitely many booleans (GX, VSTAR, Supporter, \dots) \\
History predicates & yes & booleans such as ``during your last turn'' (\textsf{Q} \cite{apr}) \\
Lingering effects & yes, if no card-recycling loop replays a modifier card within a turn & last at most about two turns, except four rest-of-game templates: three from once-per-game GX attacks and one Trainer that sets a flag; within a turn their total is bounded by the copies of each card times a printed constant, unless a card-recycling loop replays the card (Section~\ref{sec:linger-printed}) \\
Whose turn, phase & yes & finite \\
Turn number & \textbf{no} & no turn limit; card text reads it only as ``first turn'' (\textsf{I}) \\
Sudden Death games & \textbf{no} & each is a fresh game; no state carries over (\textsf{Q} \cite{rb-ssp}) \\
Coin-flip history & not state & not stored across turns (\textsf{I}) \\
\bottomrule
\end{tabularx}
\caption{State components of the printed game, after the rules audit.}\label{tab:inventory}
\end{table}

%======================================================================
\section{Survey statistics}\label{app:survey}
%======================================================================

\paragraph{Census.} Of the 20,776 card records, removing reprints left 22,745 distinct text
items; an item is one text field (an ability, an attack or a rules box) of one card name, and a
reprint is an item with the same name, field and text. These split into 36,239 sentences and
9,501 distinct templates, of which 6,963 occur once. The 141 most frequent templates cover 50\%
of sentences, 2,396 cover 80\%, and 5,878 cover 90\%. A by-eye sample put the disputed share of
the keyword-rule class labels at about 10\%.

\paragraph{Reading by eye.} The census read every distinct sentence selected by the risk-class
patterns: 170 HP-bearing templates, 68 Knock Out templates, 18 marker templates, 18 ``this
game'' templates, the 87 quantity-reading sentences whose object was not obviously bounded, 29
stacking templates, and 76 outside-the-game templates. The flip-until-tails count, 194
sentences, is the number of the 36,239 sentences that the census script's regular expression
\texttt{until you get tails|until tails|flip a coin until} matches.

\paragraph{Targeted reading.} The targeted reading classified 1,616 distinct sentences with a
duration cue, read all 143 ``as often as you like'' and ``any number of times'' sentences,
found 193 distinct sentences that mention HP, of which 65 folded templates (85 card and number
variants) change HP or maximum HP, and classified 1,654 distinct Bench sentences by regular
expression (precision spot-checked, recall not measured).

%======================================================================
\section{The Lean audit}\label{app:audit}
%======================================================================

\paragraph{Axioms.} No theorem in this paper uses an axiom other than \lean{propext} and
\lean{Quot.sound} (some use fewer). None uses \lean{Classical.choice}, and the witness for
Theorem~\ref{thm:finite} is an explicit list. The development uses core Lean only. Core's
\lean{List.perm\_iff\_count} depends on \lean{Classical.choice}, so \texttt{ListPerm.lean}
proves the needed direction (equal counts give a permutation, for a type with decidable
equality) constructively.

\paragraph{Enforcement.} \texttt{Audit.lean} runs a single command,
\lean{\#audit\_namespace Pokemon} (defined in \texttt{AuditCmd.lean}), which collects with
\lean{Lean.collectAxioms} the axioms of every theorem declared in the project's modules and in
the audit file itself, private, nested and compiler-generated theorems included, and raises a
build error if any depends on an axiom other than \lean{propext} or \lean{Quot.sound}. Since
\lean{sorry} elaborates to the axiom \lean{sorryAx}, an unfinished proof fails the same check.
The command also fails if a file under \texttt{Pokemon/} is not imported by the audit file, or
if it finds no hand-written theorem. The continuous-integration workflow builds the development
and runs \texttt{scripts/audit.sh}, which fails unless the build log contains the command's
success line, and fails on any line matching \texttt{error} or \texttt{sorry}. In tests with
planted violations (a proof by excluded middle, a \lean{sorry}, a private offender, a theorem
outside the namespace, and an unimported file), each failed the build. One gap is known: a
theorem written in the root file \texttt{Pokemon.lean}, which contains only imports, is not
checked.

\paragraph{Counts.} We built the development at commit \texttt{ce0de3b} with \texttt{lake
build} and read the audit's success line: ``\texttt{audit: 458 theorems (+1122 generated),
axioms $\subseteq$ \{propext, Quot.sound\}}''.

\subsection*{Further definitions and lemmas}\label{app:defs}

The following Lean text is cited in Sections~\ref{sec:model} and~\ref{sec:finite-side}. The
game state (docstrings omitted; Section~\ref{sec:model} describes the fields):
\begin{lstlisting}
structure GameState where
  cur : PlayerState
  opp : PlayerState
  curIsFirst : Bool
  firstTurn : Bool
  attacked : Bool
  between : Bool
  winner : Option Bool
  stadium : Option (Bool × Card)
\end{lstlisting}
Two cases of \lean{Run}
(\lean{iteF} is symmetric; \lean{selfDamage} and \lean{healAll} change the damage on the
player's Active as their comments say):
\begin{lstlisting}
  | iteT {t : Test} {a b : Effect} {g g' : GameState} {p : InPlay}
      (h : g.activeOf t.whose = some p) (ht : t.eval p = true) :
      Run a g g' → Run (.ite t a b) g g'
  | scale {g : GameState} {p q : InPlay} (w : Whose) (k : Nat)
      (hs : g.activeOf w = some p) (h : g.opp.active = some q) :
      Run (.scale w k) g
        { g with opp := { g.opp with active := some { q with damage := q.damage + k * p.damage } } }
\end{lstlisting}
The move that ends the turn:
\begin{lstlisting}
  | endTurn {g : GameState} {a b : InPlay} (fa sa ca fb sb cb : Bool)
      (ha : g.opp.active = some a) (hb : g.cur.active = some b) (hk : DamageBounded g)
      (hbt : g.between = false) :
      Move g
        { cur := { g.opp with
            active := some (a.checkup false fa sa ca),
            energyAttached := false, supporterPlayed := false, retreated := false,
            mods := g.opp.mods.shift },
          opp := { g.cur.unfresh with
            active := some (b.unfresh.checkup true fb sb cb), mods := g.cur.mods.shift,
            hadTurn := true },
          curIsFirst := !g.curIsFirst, firstTurn := false, attacked := false,
          between := true, winner := g.winner, stadium := g.stadium }
\end{lstlisting}
Lemma~\ref{thm:flag} (docstrings omitted; all in namespace \lean{Pokemon.Markers}):
\begin{lstlisting}
inductive CountOp where
  | place
  | removeAll

def CountOp.onCount : CountOp → Nat → Nat
  | .place, n => n + 1
  | .removeAll, _ => 0

def CountOp.onFlag : CountOp → Bool → Bool
  | .place, _ => true
  | .removeAll, _ => false

def hasAny (n : Nat) : Bool := Nat.blt 0 n

theorem flag_exact (ops : List CountOp) (n : Nat) :
    hasAny (ops.foldl (fun k o => o.onCount k) n) =
      ops.foldl (fun b o => o.onFlag b) (hasAny n)
\end{lstlisting}
A rest-of-game modifier survives a shift:
\begin{lstlisting}
theorem Store.add_always_shift (t : Store) (m : Mod) :
    (t.add .always m).shift = t.shift.add .always m
\end{lstlisting}
Knock Outs inside Item and Supporter effects:
\begin{lstlisting}
inductive Resolve : GameState → GameState → Prop where
  | done {g : GameState} (h : Settled g ∨ g.winner ≠ none) : Resolve g g
  | step {g g₁ g' : GameState} (hn : g.winner = none) (hm : KOMove g g₁)
      (hr : Resolve g₁ g') : Resolve g g'

inductive TRun : Effect → GameState → GameState → Prop where
  | prim {e : Effect} {g g₁ g₂ : GameState} (he : e.isPrim = true) (hn : g.winner = none)
      (hr : Run e g g₁) (hk : Resolve g₁ g₂) : TRun e g g₂
  | stop {e : Effect} {g : GameState} (hn : g.winner ≠ none) : TRun e g g
  | iteT {t : Test} {a b : Effect} {g g' : GameState} {p : InPlay} (hn : g.winner = none)
      (h : g.activeOf t.whose = some p) (ht : t.eval p = true) :
      TRun a g g' → TRun (.ite t a b) g g'
  ...
\end{lstlisting}
\noindent (The elided constructors are \lean{iteF}, symmetric to \lean{iteT}, and
\lean{seq}, \lean{choiceL}, \lean{choiceR}, \lean{iterNil} and \lean{iterCons}, which have
the same shape as in \lean{Run}.)

\noindent Theorem~\ref{thm:perm}:
\begin{lstlisting}
theorem Run.perm {e : Effect} {g g' : GameState} (h : Run e g g') :
    g'.allCards.Perm g.allCards
theorem TRun.perm {e : Effect} {g g' : GameState} (h : TRun e g g') :
    g'.allCards.Perm g.allCards
theorem Move.perm {g g' : GameState} (h : Move g g') : g'.allCards.Perm g.allCards
theorem Reach.perm {g₀ g : GameState} (h : Reach g₀ g) : g.allCards.Perm g₀.allCards
\end{lstlisting}
Theorem~\ref{thm:unbounded}:
\begin{lstlisting}
theorem counters_unbounded (n : Nat) :
    ∀ (g : GameState) (p : InPlay), g.opp.active = some p →
      ∃ g' p', Run (.iter (.counters 1)) g g' ∧ g'.opp.active = some p' ∧
        p'.damage = p.damage + n
theorem poisonUp_unbounded (n : Nat) :
    ∀ (g : GameState) (p : InPlay), g.opp.active = some p → p.poisoned = true →
      ∃ g' p', Run (.iter (.poisonUp 1)) g g' ∧ g'.opp.active = some p' ∧
        p'.poisonUp = p.poisonUp + n
\end{lstlisting}
Theorem~\ref{thm:progress}:
\begin{lstlisting}
theorem Reach.wellFormed {g₀ g : GameState} (h₀ : WellFormed g₀) (h : Reach g₀ g) :
    WellFormed g
theorem Reach.progress {g₀ g : GameState} (h₀ : WellFormed g₀) (h : Reach g₀ g)
    (hn : g.winner = none) : ∃ g', Move g g'
theorem Reach.knockOut_of_pending {g₀ g : GameState} (h₀ : WellFormed g₀) (h : Reach g₀ g)
    (hn : g.winner = none) (hk : ¬ DamageBounded g) :
    ∃ g', Move g g' ∧
      g'.cur.prizes.length + g'.opp.prizes.length + 1 =
        g.cur.prizes.length + g.opp.prizes.length
\end{lstlisting}
Theorem~\ref{thm:untouched}:
\begin{lstlisting}
theorem Move.pending_untouched {g g' : GameState} {p : InPlay} (hm : Move g g') :
    (g.cur.active = some p → p.top.hp ≤ p.damage →
        (g'.cur.active = some p ∨ g'.cur.active = none) ∧ g'.curIsFirst = g.curIsFirst) ∧
      (g.opp.active = some p → p.top.hp ≤ p.damage →
        (g'.opp.active = some p ∨ g'.opp.active = none) ∧ g'.curIsFirst = g.curIsFirst)
\end{lstlisting}
Theorem~\ref{thm:guarded}:
\begin{lstlisting}
theorem TRun.iff_guarded {e : Effect} {g g' : GameState} (hs : Settled g ∨ g.winner ≠ none) :
    TRun e g g' ↔ TRunG e g g'
theorem TRun.settled_or_won {e : Effect} {g g' : GameState} (h : TRun e g g')
    (hs : Settled g ∨ g.winner ≠ none) : Settled g' ∨ g'.winner ≠ none
\end{lstlisting}

\paragraph{Witnesses for Item timing.} \texttt{Trainer.lean} gives three witnesses, each a
single legal move from a well-formed, settled start. Here \lean{mon} is a Basic \pkmn{} with
50 HP and \lean{big} one with 100 HP.

\begin{theorem}[\lean{selfKO\_reachable}, \lean{koThenHit\_reachable},
\lean{winMidEffect\_reachable}]\label{thm:witness}
(i) An Item \lean{seq (selfDamage 5) (heal 5)} Knocks Out its player's own 50-HP Active
\pkmn{} after the first step; the opponent takes a Prize card, \lean{big} is promoted, and
the heal acts on it. (ii) An Item \lean{iter (counters 5)} Knocks Out the opponent's 50-HP
Active \pkmn{}, the opponent promotes \lean{big}, and the next round places 5 damage counters
on it. (iii) An Item \lean{seq (counters 5) draw} takes the player's last Prize card in its
first step; the game is over and the draw is not made.
\begin{lstlisting}
theorem selfKO_reachable :
    ∃ g₀, WellFormed g₀ ∧ Settled g₀ ∧ g₀.cur.active = some (InPlay.ofBasic mon) ∧
      ∃ g, Move g₀ g ∧ Reach g₀ g ∧ Card.pokemon mon ∈ g.cur.discard ∧
        g.cur.active = some (InPlay.ofBasic big) ∧
        g.opp.prizes.length + 1 = g₀.opp.prizes.length ∧ g.winner = none
theorem koThenHit_reachable :
    ∃ g₀, WellFormed g₀ ∧ Settled g₀ ∧ g₀.opp.active = some (InPlay.ofBasic mon) ∧
      InPlay.ofBasic big ∈ g₀.opp.bench ∧
      ∃ g, Move g₀ g ∧ Reach g₀ g ∧ Card.pokemon mon ∈ g.opp.discard ∧
        g.opp.active = some { InPlay.ofBasic big with damage := 5 } ∧
        g.cur.prizes.length + 1 = g₀.cur.prizes.length ∧ g.winner = none
theorem winMidEffect_reachable :
    ∃ g₀, WellFormed g₀ ∧ Settled g₀ ∧ g₀.winner = none ∧
      ∃ g, Move g₀ g ∧ Reach g₀ g ∧ g.winner = some g₀.curIsFirst ∧
        g.cur.deck = g₀.cur.deck
\end{lstlisting}
(All three are in namespace \lean{Pokemon.Trainer}.)
\end{theorem}
\noindent Witness (ii) is rulebook behaviour only under one reading of the effect. Read as ``as
often as you like, put 5 damage counters on the opponent's Active \pkmn{}'', each round is a
separate act and the Knock Out comes between rounds. Read as ``flip a coin until you get tails;
5 damage counters for each heads'', the printed card places all its damage counters at once,
and the outcome in which the second round hits the promoted \pkmn{} is a consequence of the
first caveat of Section~\ref{sec:departures}.

%======================================================================
\section{Full table of claims}\label{app:claims}
%======================================================================

Table~\ref{tab:claims} lists every substantive claim of the paper. Each status is one of the
following, or several joined by a semicolon. \emph{Machine-checked}: a Lean theorem at commit
\texttt{ce0de3b} that passes the audit of Appendix~\ref{app:audit}; ``in the model'' marks a
property of the model's definitions rather than a named theorem. \emph{False
(machine-checked counterexample)}: the claim as stated is refuted by such a theorem.
\emph{Argued}: a proof by hand, not mechanized; \emph{argued, unverified} marks a hand argument
that no second reader has checked. \emph{By definition}: true of the model's definitions as
stated. \emph{Surveyed}: found by a systematic but non-exhaustive search of rules documents or
card text. \emph{Conjectured}: no proof is offered. \emph{Open}: a question on which we take no
side. \emph{Not surveyed}: outside the survey.

{\footnotesize
\begin{xltabular}{\textwidth}{@{}>{\raggedright\arraybackslash}X>{\raggedright\arraybackslash}p{3.6cm}l@{}}
\caption{Every claim of the paper, with its status.}\label{tab:claims}\\
\toprule
Claim & Status & Where \\
\midrule
\endfirsthead
\toprule
Claim & Status & Where \\
\midrule
\endhead
\bottomrule
\endfoot
\multicolumn{3}{@{}l}{\emph{The model in general}} \\
Legal play permutes the cards in the game & machine-checked & Thm.~\ref{thm:perm} \\
After an attack or a failed Confused attempt, every move is a Knock Out move or ends the turn, giving exactly the state that ending the turn gives & machine-checked & Thm.~\ref{thm:turn} \\
Play preserves consistent turn flags; under them no attack and no Supporter on the first turn of the game, and no evolution on it & machine-checked & Thm.~\ref{thm:turn} \\
The second player's evolution ban on their own first turn & by definition (a premise of \lean{evolve}) & Thm.~\ref{thm:turn} \\
An attack with base damage 0 does no damage, whatever the modifiers & machine-checked & Sec.~\ref{sec:damage-calc} \\
A new Poison replaces the old one and resets Poison strength & machine-checked & Sec.~\ref{sec:model} \\
Typed costs are a matching of symbols to units of Energy, decided by \lean{payable} & machine-checked & Sec.~\ref{sec:damage-calc} \\
A well-formed attack makes the damage calculation at most once; Trainer effects make none & machine-checked; by definition & Sec.~\ref{sec:damage-calc} \\
With the stores capped at $k$, Resistance separates two bonuses equal at the cap (a fact about the damage function, not a game-level witness) & machine-checked & Sec.~\ref{sec:clamp} \\
Inside one effect, damage and Poison strength on the opponent's Active are unbounded & machine-checked & Thm.~\ref{thm:unbounded} \\
The unclamped reachable set from a well-formed start is infinite & machine-checked & Thm.~\ref{thm:infinite} \\
Even the reachable states with no pending Knock Out are infinite & machine-checked & Thm.~\ref{thm:infinite} \\
Every move that passes the turn yields a state with no pending Knock Out & false (machine-checked counterexample) & Rem.~\ref{rem:pendingko} \\
The turn proper begins with no pending Knock Out & machine-checked & Rem.~\ref{rem:pendingko} \\
Well-formedness is preserved; no reachable state without a winner is stuck; a pending Knock Out can be taken & machine-checked & Thm.~\ref{thm:progress} \\
A move changes an Active at or above HP only by Knocking it Out & machine-checked & Thm.~\ref{thm:untouched} \\
No step or test of an Item or Supporter effect runs while a \pkmn{} is at or above its HP & machine-checked & Thm.~\ref{thm:guarded} \\
An Item can Knock Out, promote and win in mid-effect & machine-checked & Thm.~\ref{thm:witness} \\
Knock Out timing inside Item and Supporter effects is the Advanced Player's Rulebook's, up to two caveats that lose outcomes & argued & Sec.~\ref{sec:departures} \\
No English rulebook states Knock Out checks outside attacks and Checkup & surveyed (2023 rulebook; 2024 as quoted) & Sec.~\ref{sec:game} \\
Each primitive step damages at most one Active \pkmn{}; \lean{Resolve} always has a derivation from a well-formed state & argued & Sec.~\ref{sec:finite-scope} \\
Effects read only damage, Special Conditions, markers and type of the two Actives, and the modifiers in effect through \lean{damage} and \lean{drain} & by definition & Sec.~\ref{sec:effects} \\
For place-one, remove-all and ``has any'', a flag answers as the count does & machine-checked (abstract lemma, not about \lean{Move}) & Lem.~\ref{thm:flag} \\
No card counts Imprison or Shock-wave markers or removes one at a time & surveyed (English); assumed by the model & Sec.~\ref{sec:markers} \\
Four departures are load-bearing for finiteness; the lingering modifiers enter through the tameness hypothesis & argued & Sec.~\ref{sec:departures} \\
\multicolumn{3}{@{}l}{\emph{The finite side}} \\
The clamp is a simulation inside one attack's effect & false (machine-checked counterexample) & Thm.~\ref{thm:notsim} \\
During a turn, the player's own Active has at most \lean{maxHp}${}+3$ damage counters & false (machine-checked counterexample) & Lean: \lean{Reads.activeLe\_false} \\
Tame decks: an attack's effect is simulated with own damage exact and the opponent's capped; a Trainer effect up to the clamp & machine-checked & Thm.~\ref{thm:clamp} \\
Tame decks: states with the same clamp have the same moves and reach the same clamped states & machine-checked & Thm.~\ref{thm:clamp} \\
Tame decks: the clamps of reachable states lie in one finite list, and the clamp is a two-way bisimulation along reachable states (\lean{Reach.finite}) & machine-checked & Thm.~\ref{thm:finite} \\
Every deck: the clamps of reachable states lie in one finite list (\lean{Reach.clamp\_mem}); on its own this says nothing about play & machine-checked & Thm.~\ref{thm:finite} \\
Tame decks: no fixed start simulates a two-counter machine with infinitely many reachable configurations & argued & Cor.~\ref{cor:counter} \\
Tame decks: reachability is decidable & conjectured & Rem.~\ref{rem:decide} \\
Every deck: states with the same exact-store clamp have the same moves, both ways, reachable or not (\lean{GameState.xclamp\_bisim}) & machine-checked & Thm.~\ref{thm:xbisim} \\
Every well-formed start: among reachable states with stores bounded by $B$, the exact-store clamps lie in one finite list (\lean{Reach.xfinite\_bounded}); the bounded part is not closed under moves & machine-checked & Sec.~\ref{sec:faithful} \\
Faithful starts: every reachable state's stores are bounded by \lean{storeBound} (\lean{Faithful.reach\_boundedStores}) & machine-checked & Thm.~\ref{thm:xfinite} \\
Faithful well-formed starts: the exact-store clamps of reachable states lie in one finite list, with the two-way bisimulation (\lean{Faithful.xfinite}) & machine-checked & Thm.~\ref{thm:xfinite} \\
The faithful bound rests on the model's lack of card recycling, repeated attacks and extra turns, each of which exists on printed cards & by definition (of the model); surveyed (cards) & Sec.~\ref{sec:faithful} \\
Neither the tame nor the faithful class contains the other & argued & Sec.~\ref{sec:faithful} \\
Faithful decks: reachability is decidable & conjectured & Rem.~\ref{rem:decide} \\
\multicolumn{3}{@{}l}{\emph{Where the clamp fails}} \\
For every cap $k \ge 3$ there is a deck with a damage reduction, and for every $k \ge 2$ a deck with \lean{drain}, for which \lean{clamp k} is not a simulation & machine-checked & Thm.~\ref{thm:breaks} \\
\multicolumn{3}{@{}l}{\emph{The universal side (decks that no printed card supplies)}} \\
For every two-counter machine $M$ and inputs $a$, $b$, a well-formed start from which a halt-marked state is reachable by some sequence of legal moves iff $M$ halts, every reached configuration being encoded by a reachable state (\lean{Counter.universal}) & machine-checked & Thm.~\ref{thm:universal} \\
Soundness: no sequence of legal moves reaches the halt marker unless $M$ halts & machine-checked & Thm.~\ref{thm:universal} \\
A state encodes at most one configuration & machine-checked & Thm.~\ref{thm:universal} \\
The two-counter starts satisfy \lean{LoopFree} and, when the first instruction is not \lean{halt}, violate \lean{GXOnly}, so that clause cannot be dropped (a sharpness result, not a classification of decks) & machine-checked & Thm.~\ref{thm:sharp} \\
For \lean{start countUp 0 0} the exact-store clamps of reachable states lie in no finite list (representation, not play) & machine-checked & Thm.~\ref{thm:sharp} \\
Reachability of the halt marker over the constructed starts is undecidable & machine-checked reduction; Minsky's theorem cited, not formalised & Sec.~\ref{sec:counter-what} \\
For one witness deck, no finite list covers the reachable states up to an equivalence fine enough to recover the encoded configuration (a statement about representation, not play) & machine-checked & Thm.~\ref{thm:notfiniteupto} \\
No uniform computable finite abstraction answers the halt question over the constructed starts & argued (from the reduction and Minsky's theorem) & Sec.~\ref{sec:counter-what} \\
For each single machine, a finite quotient respecting the halt question exists & argued (non-constructive) & Sec.~\ref{sec:counter-what} \\
Rest-of-game damage modifiers come only from once-per-game attacks; same-turn ones are bounded by copies times a constant unless a loop replays them & surveyed (English) & Sec.~\ref{sec:linger-printed} \\
An in-turn infinite loop has existed in a real format (Expanded) & surveyed (secondary source) & Sec.~\ref{sec:linger-printed} \\
\multicolumn{3}{@{}l}{\emph{What remains open}} \\
With attacks as the only readers, boundedly many bits cross each turn boundary & false as stated; argued, unverified & Sec.~\ref{sec:linger-next} \\
Guess-and-verify with a coupled replayable writer gives universality without the rest-of-game slot & argued, unverified & Sec.~\ref{sec:linger-next} \\
A replayable carried damage modifier exists on printed cards (Iron Defender, an Item), and it is not coupled & surveyed & Sec.~\ref{sec:linger-next} \\
Inside one effect \lean{LoopFree} excludes the guess-and-verify route; across effects the model lacks the replays it needs & argued & Sec.~\ref{sec:linger-next} \\
Finiteness once card recycling, repeated attacks and extra turns are in the model & open & Sec.~\ref{sec:linger-next} \\
For today's Standard pool: the faithful theorem once card recycling and repeated attacks are in the model & open & Sec.~\ref{sec:linger-next} \\
\multicolumn{3}{@{}l}{\emph{The printed game}} \\
Card conservation; no tokens; Bench, Prizes and deck bounded & surveyed (rulebooks) & Sec.~\ref{sec:game} \\
Damage exceeds HP only inside one attack or one Checkup & surveyed (rulebook, rulings) & Sec.~\ref{sec:game} \\
No unbounded HP, no Knock Out immunity, no game-long count, no cards from outside the game & surveyed (English) & Sec.~\ref{sec:findings} \\
Imprison and Shock-wave marker counts are uncapped but read only as ``any'' & surveyed (English) & Sec.~\ref{sec:findings} \\
No unbounded store that card text reads beyond a yes-or-no test, except possibly through card-recycling loops & surveyed (English) & Sec.~\ref{sec:survey} \\
No forced attack; Checkup is the only autonomous loop & surveyed (English) & Sec.~\ref{sec:control} \\
Damage counters can be added, healed, moved, tested for zero and for an exact number, and doubled; no halving or parity test found & surveyed & Sec.~\ref{sec:hide-stack} \\
One printed attack cannot loop over a state it reads & surveyed; argued & Sec.~\ref{sec:hide-attack} \\
Special Conditions are finite control (16 combinations per Active) & machine-checked in the model; surveyed (rules) & Sec.~\ref{sec:hide-cond} \\
Poison strength only increments, resets when Poison ends, and is read only through Checkup & machine-checked in the model; surveyed (cards) & Sec.~\ref{sec:hide-poison} \\
The printed game (English pool, current rules) is not Turing complete & conjectured & Conj.~\ref{conj:main} \\
Unbounded HP would make the game universal & conjectured & Sec.~\ref{sec:lattice} \\
Deck size as a parameter would make the game universal & conjectured & Sec.~\ref{sec:lattice} \\
Is the win probability above $p$? & open & App.~\ref{app:future} \\
Of 124 catalogued rules and mechanisms, 25 faithful, 31 simplified, 68 absent, at a time before M1 and M4 & surveyed & Sec.~\ref{sec:standard} \\
Standard (regulation marks H, I, J and Basic Energy, an assumption) has 3,402 records, 2,164 distinct cards & surveyed & Sec.~\ref{sec:standard} \\
Standard has no GX, VSTAR, rest-of-game modifiers, extra turns, markers or Lost Zone & surveyed & Sec.~\ref{sec:standard} \\
Standard has 34 carried reductions and 5 \lean{drain} attacks (so its decks are not tame), 70 recycling cards, 4 cards that attack twice, 316 \pkmn{} with Abilities (7 unlimited), 24 flip-until-tails attacks and Energy-scaled HP & surveyed & Sec.~\ref{sec:standard} \\
No Standard effect replays a same-turn damage bonus within a turn & surveyed; argued & Sec.~\ref{sec:standard} \\
Japanese-only cards; the current rulebook & not surveyed & Sec.~\ref{sec:limits} \\
\end{xltabular}
}

%======================================================================
\section{Future work}\label{app:future}
%======================================================================

\paragraph{Order.} The model grows by milestones, in the order that the Standard profile of
Section~\ref{sec:standard} recommends for proofs over today's Standard pool first:
\begin{enumerate}
\item M1, turn structure (done);
\item M4, typed Energy and costs, Weakness and Resistance, and the damage calculation (done);
\item M3, positions: the Bench as a target of effects, Knock Outs of Benched \pkmn{}, several
  Prize cards per Knock Out, and Tera;
\item MV, an effect vocabulary: moving cards between zones with card-kind filters, counts of
  hand, discard pile, Energy and Prizes, Special Conditions on one's own \pkmn{}, and recycling
  from the discard pile;
\item M8a, activated Abilities, once per turn and ``as often as you like'';
\item M5, Tool and Stadium moves, evolution and attachment on the Bench, and devolution;
\item M7 with M8b, lingering effects bound to one \pkmn{}, one turn of history, attacking twice,
  copy attacks, and static and triggered Abilities with Checkup effects;
\item M6, HP derived from the state;
\item M9, set-up and deck legality;
\item M2, the once-per-game flags for GX attacks and VSTAR Powers, which Standard lacks;
\end{enumerate}
and M10, which labels each nondeterministic branch with the player or the chance that chooses
it, at any point. MV brings card recycling into the model, and with it the open question of
Section~\ref{sec:linger-next}: the faithful bound must then be restated with the deck as the
bound on recycled resources. The profile estimates that M1, M4, M3 and MV together make about
68\% (60--70\%) of the Standard pool's distinct cards expressible, and that without MV no
combination of the other milestones exceeds 47\%. Benched Knock Outs must precede HP modifiers,
because a modifier that leaves play can lower a Benched \pkmn{}'s HP to or below its damage, and
the model can Knock Out only an Active \pkmn{}. Alongside the milestones remain a test of the
guess-and-verify route against the card text, a search of the card text for a test on Poison
strength, per-player card ownership, and the decidability of Remark~\ref{rem:decide}.

\paragraph{HP modifiers.} Of the 193 distinct sentences that mention HP, 65 folded templates
change HP or maximum HP. Their sources are Tools, Stadiums, Special Energy, Abilities, one
attack, copying (Ditto), and Trainer cards played as \pkmn{}. Every term found is a printed
constant, or a printed constant times a count of cards, Prizes or \pkmn{} in play, which gives
a structural bound. The largest total built from constant sources is 780 HP (Alolan
Exeggutor; built by hand, not proved optimal); Energy-scaled sources are bounded only by the
deck, at most 2,480 HP for Conkeldurr counted by cards. Under the rules, damage counters stay
when a bonus leaves play, so HP must be derived from the current state, never stored, and a
Knock Out check must follow every change that can lower HP. Three points are open in the card
text: two Dittos (\texttt{dp6}-27, \texttt{base3}-3) copy the maximum HP of the opposing Active
\pkmn{}, which is circular when both Active \pkmn{} are such Dittos; Heavy Boots and Rescue
Board make HP and Retreat Cost depend on each other on a \pkmn{} that holds both, a cycle whose
constructibility we have not verified; and the order of ``set'' forms (``its max HP is 300'')
relative to ``$+N$'' forms is not stated.

\paragraph{The Bench.} Effects damage, heal, evolve, devolve, switch, Knock Out, discard and
test Benched \pkmn{}, and many more act on ``1 of your \pkmn{}'' and so reach the Bench
implicitly. No sentence applies a Special Condition to a Benched \pkmn{}, but markers can be
placed on one. The Bench size is the minimum over the active limiters (as low as 3) and 5, or 8
when one of the raising cards applies (12 templates from 9 sources), and a shrink discards in a
stated order. 43 distinct sentences, all from one era, assume two Active \pkmn{} per player.

\paragraph{Importer and coverage.} The theorems reach as far as the effect language expresses
card text. We plan an importer from the census templates to \lean{Effect} and a checked count of
the card records whose text it translates, with the remainder named and a reason for each. The
census suggests the size of the task: 141 templates cover half of all sentences, but the last
tenth of the sentences needs 3,623 more templates beyond the 5,878 that cover 90\%. Japanese-only
cards need a corpus of their own, and the marker finding is the part most at risk.

\paragraph{Other questions.} With deck size, Bench size, HP or card pool as parameters, the
game becomes a family of finite games whose complexity is open; constraint logic
\cite{hearn-demaine-ncl,hearn-demaine-book} is a natural tool for hardness, and we have not
tried it. Finiteness does not settle quantitative questions: for a probabilistic finite
automaton the acceptance probability of a given word is computable, but the threshold question
is undecidable \cite{nasu-honda,paz,condon-lipton}, even with the numbers of states and letters
fixed \cite{blondel-canterini} and for one fixed automaton when only the starting distribution
varies \cite{rote-pfa}. We pose as open whether, for a family of positions, the probability
that the first player wins exceeds a given $p$; the question is compatible with
Theorem~\ref{thm:finite} and not answered by it, and whether coin flips and shuffles realize
the automata the undecidability proofs need is untested. Copy attacks that call copy attacks
(Togetic, Metronome, Foul Play) are a possible call stack inside effect resolution; the
control-flow search found the depth bounded by the board, and mutual copying is unresolved
(Section~\ref{sec:hide-attack}).

\end{document}
